% ====================================================================
% graphite_ica_ch1.tex — Chapter 1
%   흑연 음극 dQ/dV(ICA) 상전이 피크의 물리:
%   피크 {위치·너비·높이}가 {온도·C-rate·극판 전위}에 어떻게 반응하는가를
%   열역학·동역학·전기화학·통계로 유도해 실무 피팅에 쓸 닫힌 식에 도달한다.
% ====================================================================
\documentclass[11pt,a4paper]{article}
\usepackage[margin=22mm]{geometry}
\usepackage{kotex}
\usepackage{amsmath,amssymb,mathtools,bm}
\usepackage{booktabs,longtable,array}
\usepackage{enumitem}
\usepackage{amsthm}
\usepackage{hyperref}
\usepackage{fancyhdr}
\usepackage{lastpage}
\usepackage{setspace}
\setstretch{1.12}
\setlength{\parskip}{0.45em}\setlength{\parindent}{0pt}\setlength{\headheight}{14pt}
\hypersetup{colorlinks=true, linkcolor=blue, urlcolor=blue, citecolor=blue,
  pdftitle={흑연 음극 dQ/dV 이론: 평형·동역학·충방전 히스테리시스 — Chapter 1}}
\pagestyle{fancy}\fancyhf{}
\lhead{흑연 음극 $dQ/dV$ 이론 — Ch.1}\rhead{\thepage/\pageref{LastPage}}
\renewcommand{\headrulewidth}{0.4pt}

\newtheorem*{keybox}{핵심}
\newtheorem*{stagebox}{흑연 staging 배경}
\newtheorem*{boundbox}{유효 범위}
\newtheorem*{flagbox}{비고}
\newtheorem*{rigorbox}{심화}
\newtheorem*{verifybox}{검산}

\newcommand{\dd}{\mathrm{d}}
\newcommand{\eff}{\mathrm{eff}}
\newcommand{\eq}{\mathrm{eq}}
\newcommand{\app}{\mathrm{app}}
\newcommand{\drive}{\mathrm{drive}}
\newcommand{\cell}{\mathrm{cell}}
\newcommand{\bg}{\mathrm{bg}}
\newcommand{\ext}{\mathrm{ext}}
\newcommand{\OCV}{\mathrm{OCV}}
\newcommand{\sstat}{\mathrm{ss}}
\newcommand{\peak}{\mathrm{peak}}
\newcommand{\dis}{\mathrm{dis}}
\newcommand{\chg}{\mathrm{chg}}
\newcommand{\hys}{\mathrm{hys}}
\newcommand{\pol}{\mathrm{pol}}
\newcommand{\kin}{\mathrm{kin}}
\newcommand{\obs}{\mathrm{obs}}

\renewcommand{\theequation}{1.\arabic{equation}}
\renewcommand{\thesection}{1.\arabic{section}}
\title{\textbf{\LARGE Chapter 1}\\[0.35em]\Large 흑연 음극 $dQ/dV$ 이론: 평형 peak $\cdot$ 동역학 $\cdot$ 충방전 히스테리시스}
\author{}
\date{}

\begin{document}
\maketitle

% ====================================================================
\section*{서론}
\addcontentsline{toc}{section}{서론}
% ====================================================================
리튬이온전지 흑연 음극의 \textbf{incremental capacity analysis}(ICA)는 충방전 곡선 $Q(V)$ 의 미분
$\dd Q/\dd V$ 를 본다 \cite{bloom2005,dubarry2012}. 흑연은 리튬과 단계적 staging 상전이(stage $4\!\to\!3\!\to\!2\mathrm L\!\to\!2\!\to\!1$;
$2\mathrm L$ 은 stage-2 의 희박(dilute) 단계)를 거치므로 \cite{dahn1991,ohzuku1993,funabiki1999}, 각 상전이가
일어나는 전위에서 $\dd Q/\dd V$ 에 \textbf{peak} 가 선다. 관측되는 유효 전이는 $j=1,\dots,N_p$ 개이며
($N_p\approx3\sim4$; 인접 전이의 중심 전위 차가 전이폭 이하이면 peak 가 융합한다), \emph{각 전이마다} 전환율과
peak 가 따로 있다.

본 장이 설명하려는 \textbf{관측}은 다음이다.
\begin{quote}\itshape
$\dd Q/\dd V$ 의 상전이 peak 는 그 \textbf{위치$\cdot$너비$\cdot$높이}가 \textbf{(a) 온도, (b) C-rate(전류),
(c) 극판 전위} 세 가지에 따라 변한다. 특히 \emph{온도가 낮을수록 peak 의 꼬리(tail)가 길게 늘어져 다음 peak
와 겹치고}, 높을수록 짧게 끝난다.
\end{quote}

\textbf{이 장의 목표}는 이 관측을 화학열역학$\cdot$전기화학$\cdot$반응속도론$\cdot$통계와 미적분의 물리로,
\textbf{중간 단계를 빠짐없이 밝혀} 설명하는 \emph{닫힌 논리식}을 세우는 것이다.
첫째 목표는 그 식이 \emph{실데이터 피팅에 바로 쓸 수 있을 것}, 둘째는 그 유도가 \emph{1차 원리에서 빈틈없이 닫힐 것}이다.

\textbf{도착점은 하나의 식이다.} 본 장은 마지막 절(\S\ref{sec:master})에서 $\dd Q/\dd V$ 를 \emph{한 줄의 피팅식}으로
모으고, 어느 측정(저율 OCV$\cdot$전류 차단$\cdot$다율$\cdot$다온도)이 어느 파라미터를 닫는지까지 매듭짓는다. 그에
이르도록 \emph{각 절은 그 식의 한 항을 세운다} — 평형 상승부(\S\ref{sec:eqpeak}), 전위가 낮추는 동역학 꼬리
(\S\ref{sec:lag}$\cdot$\S\ref{sec:barrier}), 입자 분포가 만드는 늘어짐(\S\ref{sec:dist}), 여러 전이의 합산
(\S\ref{sec:overlap}). 독자는 매 절 끝에서 ``이 결과가 최종식의 무엇이 되는가''를 확인하게 된다.

\begin{stagebox}
흑연의 리튬화는 층간 갤러리를 단번에 채우지 않고 \emph{stage} 순서 $4\to3\to2\mathrm L\to2\to1$ 로 단계적으로
채운다($2\mathrm L$=stage-2 의 희박(dilute) 단계) \cite{dahn1991,ohzuku1993}. stage 번호 $n$ 은 대략 ``리튬으로
채운 한 층마다 빈 갤러리 $n-1$ 개''를 뜻하며, 채울수록 $n$ 이 줄어 stage 1$=\mathrm{LiC_6}$(완전 채움)$\cdot$
stage 2$=\mathrm{LiC_{12}}$ 가 된다(stage 3$\cdot$4$\cdot$2L 은 더 희박; 조성은 근사). \textbf{각 stage 전이가
$\dd Q/\dd V$ 에 peak 를 하나씩 남기며}, 인접 전이의 중심 전위 차가 전이폭 이하이면 융합해 관측 peak 수가
$N_p\approx3\sim4$ 가 된다 — 이것이 본 장이 전이 $j=1,\dots,N_p$ 로 다루는 대상이다. (위 채움 순서는 \emph{리튬화(충전)} 방향이며, 본 장은
\emph{방전}=탈리튬화 기준이라 진행은 그 역순 — $\theta_j$ 가 $1\to0$ 으로 줄고 stage 는 역방향으로 풀린다.)
\end{stagebox}

\tableofcontents
\newpage

% ====================================================================
\section{기호와 규약 (Notation \& Conventions)}\label{sec:notation}
% ====================================================================
본문에 들어가기 전에 기호$\cdot$단위$\cdot$부호 규약을 정리한다(리뷰$\cdot$시뮬레이션 논문 관행). 각 양은 처음
쓰이는 곳에서 다시 물리적으로 동기 부여하되, 여기 표를 사전으로 둔다.

\textbf{전이 인덱스 규약.} 상전이는 $j=1,\dots,N_p$ 로 여럿이다. \emph{각 전이의 Li 점유율은 $\theta_j$, 전환율(진행률)은 $\xi_j\equiv1-\theta_j$}(방전=탈리튬화 진행 시 $\theta_j$ 가 $1\to0$ 으로 줄어 $\xi_j$ 가 $0\to1$), 평형값
$\xi_{\eq,j}$, 속도상수 $k_j$, 정/역 방향 속도 $r_j^\pm$, 지연 $r_j\equiv\xi_{\eq,j}-\xi_j$, 꼬리 길이
$L_{q,j}$. 한 전이만 논하는 문맥에서는 $j$ 를 생략할 수 있으나, \emph{박스로 강조된 최종 결과식}에는 항상 $j$ 를
단다.

\textbf{부호 규약.} $s\in\{+1,-1\}$ = 전이의 방전 방향 부호(방전 기준 $s=+1$). $s_I\in\{+1,-1\}$ = 전류 부호
규약(방전 시 분극이 apparent 전위를 올리는 방향을 $+$ 로; 본 장 방전 기준 $s_I=+1$).

\begin{longtable}{p{0.15\textwidth} p{0.13\textwidth} p{0.64\textwidth}}
\toprule 기호 & 단위 & 의미 \\ \midrule \endhead
$\theta_j$ & --- & 전이 $j$ 의 Li 점유율(분율), 방전 시 $1\to0$ \\
$\xi_j,\ \xi_{\eq,j}$ & --- & 전이 $j$ 의 전환율(진행률 $=1-\theta_j$)과 그 평형값, $0\le\xi_j\le1$ \\
$Q_j$ & C & 전이 $j$ 가 완전 진행 시 담는 용량 ($>0$) \\
$Q_\cell$ & C & 셀 기준 총 방전 용량 \\
$Q_\bg(V_n,T)$ & C & staging 으로 분리되지 않는 배경 누적 용량 \\
$C_\bg$ & C/V & 배경 미분용량 $=\partial Q_\bg/\partial V_n$ (chemical capacitance) \\
$q$ & --- & 무차원 용량 좌표 $=Q_\ext/Q_\cell=\int|I|\,\dd t/Q_\cell$ \\
$V_n$ & V & \emph{내부} 음극 전위 — 전하 보존식의 해 \\
$V_\app$ & V & 측정되는 apparent 전위 $=V_n+s_I|I|R_n$ \\
$V_\drive$ & V & 상변이 속도를 구동하는 전위; 기본 $V_\drive=V_n$ \\
$V_{n,\OCV}$ & V & 평형($\xi_j=\xi_{\eq,j}$)에서의 전하 보존식 해 \\
$U_j(T),\ U_j^0$ & V & 전이 $j$ 의 평형 중심 전위, 그 표준값 \\
$w_j$ & V & 평형 전이 폭(이상 극한서 $=RT/F$; 일반은 effective width) \\
$R_n$ & $\Omega$ & 음극 유효 분극 저항 \\
$\mu_\mathrm{Li},\ \mu^0$ & J/mol & 삽입 리튬의 \emph{점유율} 화학퍼텐셜과 그 표준값 \\
$\Omega_j$ & J/mol & 정규용액 상호작용 에너지 척도 \\
$\mathcal A_j$ & J/mol & 전이 $j$ 의 구동력(affinity) $=sF(V_\drive-U_j)$ \\
$\Delta G_{a,j}$ & J/mol & 활성화 자유에너지 $=\Delta H_{a,j}-T\Delta S_{a,j}$ \\
$\Delta H_{a,j},\Delta S_{a,j}$ & {\footnotesize J/mol;\,J/(mol\,K)} & 활성화 엔탈피$\cdot$엔트로피(\emph{속도}에 들어감) \\
$\Delta G_{\eff,j}$ & J/mol & 유효 활성화 자유에너지 $=\Delta G_{a,j}-\chi\mathcal A_j$ \\
$\Delta S_j$ & J/(mol\,K) & 전이 $j$ 의 \emph{삽입(리튬화) 반응} 엔트로피($\partial U_j/\partial T=\Delta S_j/(sF)$; 활성화 $\Delta S_{a,j}$ 와 \emph{다른 양}) \\
$k_j,\ k_0$ & 1/s & 완화 속도상수 $=r_j^++r_j^-$, Eyring prefactor $k_0=k_BT/h$ \\
$k_j^{\eff},k_{j,\lim},k_{j,\mathrm{act}}$ & 1/s & 병목 포함 유효 속도 $1/k_j^{\eff}=1/k_{j,\mathrm{act}}+1/k_{j,\lim}$ ($k_{j,\mathrm{act}}\!\equiv\!k_j$=활성화 완화속도, $k_{j,\lim}$=직렬 율속) \\
$r_j^\pm$ & 1/s & 전이 $j$ 의 정/역 방향 속도 \\
$\xi_{\sstat,j},\ \bar r_j$ & --- & 정상점 $=r_j^+/(r_j^++r_j^-)$; 가중 평균 지연 $=\int\rho_G\,r_j\,\dd G$ \\
$\chi$ & --- & 전달 계수(transfer coefficient) $0\le\chi\le1$ (전기화학 표준 표기 $\alpha$; 본 장 $\chi$) \\
$\lambda$ & J/mol & Marcus 재구성 에너지 \\
$\rho_G(G),\ \bar G$ & mol/J;\,J/mol & 활성화 배리어 분포(용량 가중, $\int\rho_G\,\dd G=1$)와 우세 집단 대표값 \\
$\sigma_G$ & J/mol & 배리어 분포의 표준편차 \\
$q_a,\ V_a$ & ---,\ V & 꼬리 시작 좌표와 그 전위 $V_a=V_n(q_a)$ \\
$r_{a,j}$ & --- & 꼬리 시작점 잔류 지연 $=r_j(q_a)=\xi_{\eq,j}(q_a)-\xi_j(q_a)$, $0<r_{a,j}\le1$(꼬리가 전이 전부면 $=1$) \\
$L_{q,j}$ & --- & 용량축 꼬리 특성 길이 $=|I|/(Q_\cell k_j)$ (공율속 시 $k_j\!\to\!k_j^{\eff}$; \S\ref{sec:barrier}) \\
$L_{V,j}$ & V & 전압축 꼬리 길이 $=|\dd V_n/\dd q|_{q_a}L_{q,j}$ \\
$R,F,k_B,h,N_A$ & {\footnotesize J/(mol\,K), C/mol, J/K, J\,s, mol$^{-1}$} & 기체상수, Faraday, Boltzmann, Planck, Avogadro ($R=k_BN_A$) \\
\bottomrule
\end{longtable}

% ====================================================================
\section{전하 보존이 정하는 내부 전위 $V_n$}\label{sec:stage}
% ====================================================================
상전이를 구동하는 것은 음극 \emph{내부}의 전위 $V_n$ 이다. 그런데 회로에서 읽는 전압은
$V_n$ 이 아니다(분극이 섞여 있다). \emph{그렇다면 $V_n$ 은 어디서 오는가?} 그것은 외부에서 주어지는 입력이 아니라
\textbf{전하 보존이라는 한 등식의 해}다. 이 절은 그 전위 $V_n$ 의 정의를 세운다 — 이후 ``전위가 배리어를
낮춘다''(\S\ref{sec:barrier})가 그 위에서 전개된다.

\textbf{전하 보존식.} 외부 회로로 흘린 방전 전하 $Q_\cell q$ 는 음극의 \emph{상태 변화}로 정확히 \emph{계상}된다(전하를 ``저장''하는 게
아니라 — 방전은 탈리튬화로 Li 를 \emph{빼낸다} — 전하 \emph{수지}다). 그 상태는 두 몫으로 나뉜다:
(i) 뚜렷한 staging 전이로 분리되는 몫 $\sum_j Q_j\xi_j$ (전이 $j$ 가 진행률 $\xi_j=1-\theta_j$ 만큼 진행하면 — 점유율
$\theta_j$ 가 그만큼 줄면 — 그 용량 기여는 $Q_j\xi_j$), (ii) 그렇게 분리되지 않는 배경 몫 $Q_\bg(V_n,T)$. 외부로 흘린 전하 $Q_\cell q$ 가 이 두 몫의 합과
같다는 \emph{전하 수지(보존)}가
\begin{equation}
\boxed{\;Q_\cell\,q=Q_\bg(V_n,T)+\sum_{j=1}^{N_p}Q_j\,\xi_j\;.}
\label{eq:charge}
\end{equation}
좌변은 외부 입력(아는 값), 우변은 내부 상태($V_n$ 과 $\{\xi_j\}$ 의 함수)다. 주어진 $(q,T,\{\xi_j\})$ 에서 이
식을 만족하는 $V_n$ 이 곧 내부 전위다 — \emph{외부 lookup 이 아니라 보존식의 해}다(다공성 전극의 implicit
화학퍼텐셜 관계와 같은 구조 \cite{mckinnon1983,newman2004}).

\textbf{해의 유일성(배경 미분용량).} 식~\eqref{eq:charge} 가 $V_n$ 에 대해 유일해를 가지려면 우변이 $V_n$ 의
단조 함수여야 한다. 배경 용량의 기울기
\begin{equation}
C_\bg\equiv\frac{\partial Q_\bg}{\partial V_n}\ \ge\ \epsilon_Q\ >\ 0
\label{eq:cbg}
\end{equation}
가 양이면 충분하다($C_\bg$=잔류 chemical capacitance). $\epsilon_Q$(단위 C/V)는 임의 정규화 상수가 아니라 ``양의
미분용량''이라는 물리 조건이며, 단조성과 우변 목표가 $Q_\bg$ 치역 안에 있다는 조건이 함께면 중간값 정리로 해가
유일하다 \cite{bazant2013}.

\textbf{평형 극한 = OCV.} 평형(또는 충분히 느린 조건)에서는 $\xi_j=\xi_{\eq,j}(V_n,T)$(다음 절에서 정의)이고,
식~\eqref{eq:charge} 의 해가 곧 \emph{평형 음극 전위} $V_{n,\OCV}(q,T)$ 다(평형 분기에서는 우변이 $C_\bg+\sum_j
Q_j\,\dd\xi_{\eq,j}/\dd V_n$ 로 고정-$\{\xi_j\}$ 보다 더 가파른 단조라 해가 \emph{a fortiori} 유일하다 —
$\dd\xi_{\eq,j}/\dd V_n>0$ 는 다음 절). \textbf{즉 OCV 곡선은 임의의
lookup 이 아니라 전하 보존식의 평형 특수해다.} 전류가 흐르면 $\xi_j$ 가 평형을 못 따라가(다음다음 절), 그
차이를 메우려 $V_n$ 이 OCV 에서 이탈한다 — 이 이탈이 꼬리가 생기는 바탕이다.

\textbf{측정 전위와 구동 전위.} \emph{측정}되는 전위는 내부 전위에 분극(저항 강하)이 더해진 값이다:
\begin{equation}
V_\app=V_n+s_I|I|R_n.
\label{eq:vapp}
\end{equation}
상전이 \emph{속도}를 구동하는 것은 그중 구동 전위 $V_\drive$ 이며, 가장 보수적인 기본형은 $V_\drive=V_n$ 이다
(계면 과전압을 따로 분리하는 일반형은 이후 장). 세 전위 $V_n$(내부 해)$\cdot V_\app$(측정)$\cdot V_\drive$(구동)는
\emph{물리가 다른 양}이라 구분한다 — 표기 규약이 아니라 서로 다른 측정/역할이다.

\begin{keybox}
\textbf{전하 보존의 요지.} $V_n$ 은 전하 보존식~\eqref{eq:charge} 의 해다(외부값 아님). 평형이면 그 해가 OCV,
전류가 흐르면 $\xi_j$ 의 지연만큼 $V_n$ 이 OCV 에서 밀린다. 측정은 $V_\app=V_n+s_I|I|R_n$, 속도 구동은
$V_\drive=V_n$. 이 바탕 위에서 다음 절들이 $\xi_{\eq,j}$(목표)$\to$지연$\to$전위 배리어 낮춤을 차례로 전개한다.
\end{keybox}

% ====================================================================
\section{평형에서의 peak — 기준선}\label{sec:eqpeak}
% ====================================================================
전류가 거의 흐르지 않는 \emph{평형}에서, 전이 $j$ 의 peak 는 어디에 서고 얼마나 넓고 높은가? 그리고
그것이 온도에 어떻게 변하는가? 이 기준선을 먼저 세워야, 전류가 흐를 때의 변화를 그로부터의 \emph{이탈}로 읽을
수 있다.

\subsection{$\dd Q/\dd V_n$ 의 peak 발생}
방전 궤적에서 외부 전하 $Q\equiv Q_\ext=Q_\cell q$ 와 내부 상태 $\{\xi_j\}$ 가 \emph{모두} $V_n$ 을 따라 변하므로(전하
보존식~\eqref{eq:charge} 가 그 관계를 묶는다 — 평형에서는 $\xi_j=\xi_{\eq,j}(V_n)$, 아래 소절), 식~\eqref{eq:charge}
양변을 $V_n$ 으로 미분하면(등온) 관측 ICA $\dd Q/\dd V_n=\dd(Q_\cell q)/\dd V_n$ 이 닫힌다:
\begin{equation}
\frac{\dd Q}{\dd V_n}=\frac{\partial Q_\bg}{\partial V_n}+\sum_j Q_j\frac{\dd\xi_j}{\dd V_n}
=C_\bg+\sum_j Q_j\frac{\dd\xi_j}{\dd V_n}.
\label{eq:dQdV}
\end{equation}
배경 $C_\bg$ 는 완만한 기저선이고, 한 전이 $j$ 가 진행되는 좁은 전위 구간에서는 그 항 $Q_j\,\dd\xi_j/\dd V_n$ 이
급증해 \textbf{peak} 를 만든다. \emph{즉 peak 는 $\xi_j$ 가 $V_n$ 에 대해 가장 급히 변하는 곳}이다. 그러므로
peak 의 모양은 곧 $\xi_j(V_n)$ 의 기울기 모양이며, 평형에서는 $\xi_j=\xi_{\eq,j}(V_n)$ 다. 이제 $\xi_{\eq,j}$ 를
격자기체(lattice-gas) 모형으로 유도한다.

\subsection{평형 전환율 $\xi_{\eq,j}(V_n)$ — 격자기체(lattice-gas) 유도}
\textbf{유도의 출발점.} 한 전이를 ``격자 자리(site)에 리튬이 차고 비는'' 격자기체(lattice-gas)로 보면, \emph{점유율}
$\theta_j$(Li 가 자리를 채운 분율)의 평형값은 \emph{자리를 더 채우려는 화학적 경향과 전극 전위의 균형}에서 정해진다.
그 경향을 \emph{점유율의} 화학퍼텐셜 $\mu_\mathrm{Li}(\theta_j)$ 로 적어 두 출처(섞임 엔트로피, 입자 간 상호작용)를 한
줄씩 유도한 뒤, \emph{방전 진행률} $\xi_j=1-\theta_j$ 로 옮긴다 \cite{mckinnon1983,bazant2013}.

\emph{(i) 섞임 엔트로피.} $N$ 개 동등한 자리 중 $n$ 개가 점유된 배열의 수는 $W=\dfrac{N!}{n!\,(N-n)!}$ 이다.
Boltzmann 관계 $S_{\mathrm{mix}}=k_B\ln W$ 에 Stirling 근사 $\ln m!\simeq m\ln m-m$ 를 세 계승에 적용하면
\begin{align}
\ln W&=\ln N!-\ln n!-\ln(N-n)!\notag\\
&\simeq\big(N\ln N-N\big)-\big(n\ln n-n\big)-\big((N-n)\ln(N-n)-(N-n)\big)\notag\\
&=N\ln N-n\ln n-(N-n)\ln(N-n)\qquad\big[\text{선형항}\ -N+n+(N-n)=0\big].
\label{eq:lnW}
\end{align}
점유율 $\theta_j\equiv n/N$ 을 넣으면($n=N\theta_j$, $N-n=N(1-\theta_j)$, 그리고 $N\ln N$ 을 두 항에 분배하면
$N\ln N=n\ln N+(N-n)\ln N$ 이라 $\ln(n/N),\ln((N-n)/N)$ 로 묶인다)
\begin{equation}
S_{\mathrm{mix}}=-k_B N\big[\theta_j\ln\theta_j+(1-\theta_j)\ln(1-\theta_j)\big].
\label{eq:smix}
\end{equation}
혼합 자유에너지는 $G_{\mathrm{mix}}=-T\,S_{\mathrm{mix}}$ 다. 이를 \emph{격자 자리(전체 site) 1몰당}으로 환산하면(전체
자리 수 $N$ 으로 나누고 $N_A$ 를 곱해 $R=k_B N_A$), 1몰당 혼합 자유에너지 $\bar G_{\mathrm{mix}}=RT[\theta_j\ln\theta_j+
(1-\theta_j)\ln(1-\theta_j)]$ 이고, 화학퍼텐셜의 섞임 몫은 이를 $\theta_j$ 로 미분한 것이다:
\begin{equation}
\frac{\partial\bar G_{\mathrm{mix}}}{\partial\theta_j}=RT\big[(\ln\theta_j+1)-(\ln(1-\theta_j)+1)\big]
=RT\ln\frac{\theta_j}{1-\theta_j}
\label{eq:mumix}
\end{equation}
($\theta_j\ln\theta_j$ 미분 $=\ln\theta_j+1$, $(1-\theta_j)\ln(1-\theta_j)$ 미분 $=-(\ln(1-\theta_j)+1)$, 두 $+1$ 이 상쇄).

\emph{(ii) 상호작용.} 정규용액(regular-solution) 근사는 \emph{평균장}에서 점유-인접 상호작용을 본다 — 무작위
배치에서 ``점유$\cdot$빈자리'' 쌍이 생길 확률이 $\propto\theta_j(1-\theta_j)$ 이므로 혼합 엔탈피가
$\bar G_{\mathrm{int}}=\Omega_j\,\theta_j(1-\theta_j)$ 이다 \cite{mckinnon1983}. 미분하면(곱미분
$\dd[\theta_j-\theta_j^2]/\dd\theta_j=1-2\theta_j$)
\begin{equation}
\frac{\partial\bar G_{\mathrm{int}}}{\partial\theta_j}=\Omega_j(1-2\theta_j).
\label{eq:muint}
\end{equation}
표준 상태 $\mu^0$ 까지 더해 삽입 리튬의 \emph{점유율} 화학퍼텐셜은
\begin{equation}
\mu_\mathrm{Li}(\theta_j)=\mu^0+RT\ln\frac{\theta_j}{1-\theta_j}+\Omega_j(1-2\theta_j).
\label{eq:mu}
\end{equation}

\emph{(iii) 전기화학 평형 조건 — 점유율 균형에서.} 식~\eqref{eq:mu} 의 $\mu_\mathrm{Li}$ 는 \emph{점유율
$\theta_j$ 의} 화학퍼텐셜이다(점유율이 높을수록 큼). 전극에서의 삽입 반응 $\mathrm{Li}^++e^-+[\,]\rightleftharpoons
\mathrm{Li}_{(\text{graphite})}$ ($[\,]$=흑연 빈 자리; 각 stage 전이 $j$ 를 이 반응의 \emph{유효}(coarse-grained)
형으로 본다)의 평형은 양변의 전기화학퍼텐셜이 같은 조건이다: $\mu_\mathrm{Li}(\theta_{\eq,j})=\tilde\mu_{\mathrm{Li}^+}
+\tilde\mu_{e^-}$. 전자의 전기화학퍼텐셜은 $\tilde\mu_{e^-}=\mu_{e^-}^0-FV_n$ 로 \emph{전기적 일} 항 $-FV_n$ 을
갖는다(전자 1몰을 전위 $V_n$ 전극에 놓는 일; 전하 $-F$ 가 전위 $V_n$ 과 곱). 전해질 활동도가 일정한 조건에서
$\mathrm{Li}^+\cdot$표준 항을 한 상수로 묶어 $\tilde\mu_{\mathrm{Li}^+}+\mu_{e^-}^0\equiv sF\,U_j^0$ 로 전위 $U_j^0$ 를
정의하면, 평형은 \emph{점유율} 화학퍼텐셜이 전위에 묶인 형이 된다(상수 $\mu^0$ 을 $U_j$ 로 흡수; 전기적 일 $-FV_n$ 은
$s$-무관이나 위 정의가 $s$ 를 $U_j^0$ 에 흡수하므로 $V_n$ 항이 $sF$ 로 묶이며, $s=+1$ 에서 $-F$ 와 일치한다):
\begin{equation}
\mu_\mathrm{Li}(\theta_{\eq,j})=\mu^0-sF\big(V_n-U_j\big),\qquad U_j\equiv U_j^0-\frac{\mu^0}{sF}
\label{eq:eqcond}
\end{equation}
\textbf{부호의 물리적 의미.} 좌변 $\mu_\mathrm{Li}$ 는 점유율이 높을수록 크고, $V_n$ 이 오르면(전자 일 $-FV_n$) 평형
점유율이 \emph{낮아져야} 하므로 우변이 $V_n$ 에 \emph{감소}한다 — 이 음($-$)부호가 곧 ``$V_n\uparrow\Rightarrow$
탈리튬화 진행''의 물리다(평형을 진행률 $\xi_j$ 가 아니라 점유율 $\theta_j$ 로 적었으므로 이 부호가 $s$ 규약에 가려지지 않고 명시로
드러난다). 부호 $s$ 는 진행 방향 규약이며 그 값은 \emph{관측으로 고정}된다 — 방전에서 $V_n$ 이 오를수록 진행률
$\xi_j=1-\theta_j$ 가 \emph{커지므로} $\dd\xi_{\eq,j}/\dd V_n>0$(식~\eqref{eq:dxidV})이고, 따라서 $s=+1$ 이다(충전
branch 는 $s=-1$; 다음 장).

\emph{(iv) 진행률 $\xi_{\eq,j}$ 로 옮겨 풀기.} \textbf{본 장 본문은 이상 극한 $\Omega_j=0$ 을 채택한다}(상호작용
보정 $\Omega_j\ne0$ 의 효과는 아래 비고와 \S\ref{sec:falsify} 형태 판별로 미룬다). 폭을 $w_j\equiv RT/F$ 로
두면 식~\eqref{eq:eqcond}($\Omega_j=0$)은 $RT\ln[\theta_{\eq,j}/(1-\theta_{\eq,j})]=-sF(V_n-U_j)$ 다. 이제
\emph{점유율을 방전 진행률로} 옮긴다 — $\theta_{\eq,j}=1-\xi_{\eq,j}$ 이므로 $\ln\dfrac{\theta_{\eq,j}}{1-\theta_{\eq,j}}
=\ln\dfrac{1-\xi_{\eq,j}}{\xi_{\eq,j}}=-\ln\dfrac{\xi_{\eq,j}}{1-\xi_{\eq,j}}$ 로 \emph{부호가 한 번 명시적으로
뒤집힌다}(이 반전이 ``점유율$\downarrow$ = 진행률$\uparrow$''의 대수적 표현이다):
\begin{align}
-RT\ln\frac{\xi_{\eq,j}}{1-\xi_{\eq,j}}&=-sF(V_n-U_j)\notag\\
\ln\frac{\xi_{\eq,j}}{1-\xi_{\eq,j}}&=\frac{s(V_n-U_j)}{w_j}\qquad[\div(-RT),\ w_j=RT/F]\notag\\
\frac{\xi_{\eq,j}}{1-\xi_{\eq,j}}&=E,\qquad E\equiv\exp\!\Big[\frac{s(V_n-U_j)}{w_j}\Big]\qquad[\text{양변 지수}]\notag\\
\xi_{\eq,j}&=E(1-\xi_{\eq,j})\ \Rightarrow\ \xi_{\eq,j}(1+E)=E\ \Rightarrow\ \xi_{\eq,j}=\frac{E}{1+E}.\notag
\end{align}
분자$\cdot$분모를 $E$ 로 나누면 $1/E=\exp[-s(V_n-U_j)/w_j]$ 이므로
\begin{equation}
\boxed{\;\xi_{\eq,j}(V_n,T)=\frac{1}{1+\exp\!\big[-s\,(V_n-U_j(T))/w_j(T)\big]}\;,\qquad w_j=\frac{RT}{F}\;.}
\label{eq:logistic}
\end{equation}
이것은 \emph{매끄러운} logistic 곡선이지 $0\!\to\!1$ 급점프가 아니다(본 장은 어떤 step-function 도 쓰지 않는다).

\begin{flagbox}
\textbf{평형 \emph{형태}는 데이터가 정한다.} 위 logistic 은 이상 lattice-gas($\Omega_j=0$)의 닫힌 결과다.
$\Omega_j\ne0$ 이면 식~\eqref{eq:eqcond} 에 상호작용 항 $\Omega_j(1-2\theta_{\eq,j})$(진행률로는 $-\Omega_j(1-2\xi_{\eq,j})$)가 살아 \emph{음함수}가 되어 단순
logistic 이 아니며, $w_j$ 를 effective width 로 두는 것은 닫힌 유도가 아니라 coarse-grained 근사다. 특히
$\Omega_j>2RT$ 면 isotherm 이 비단조가 되어 \emph{상분리}(2-상 공존, miscibility gap)가 나타나고, 그 구간 peak
는 본 단일-phase logistic 로 기술되지 않는다(흑연 staging 의 핵심 물리 — 본 장 범위 밖). ``Gaussian
(erf) 형''으로 보는 것도 \emph{유추}일 뿐 확정이 아니다. 본 장의 꼬리 \emph{원거리 감쇠 길이$\cdot$부호} 결론은
평형 형태 세부에 둔감하다 — 포화 후 꼬리에서는 $\xi_{\eq,j}$ 가 이미 포화해 $\dd\xi_{\eq,j}/\dd V_n\simeq0$ 이기
때문이다(꼬리 \emph{시작 위치$\cdot$진폭}은 peak 근처 평형 형태에 일부 의존). 형태 판별은 \S\ref{sec:falsify}.
\end{flagbox}

\subsection{평형 peak 의 위치$\cdot$너비$\cdot$높이}
\textbf{도함수.} $z\equiv s(V_n-U_j)/w_j$ 로 두면 $\xi_{\eq,j}=(1+e^{-z})^{-1}$ 이다. 직접 미분하면
\begin{equation}
\frac{\dd\xi_{\eq,j}}{\dd z}=-\,(1+e^{-z})^{-2}\cdot(-e^{-z})=\frac{e^{-z}}{(1+e^{-z})^2}.
\label{eq:dxidz}
\end{equation}
한편 $1-\xi_{\eq,j}=\dfrac{e^{-z}}{1+e^{-z}}$ 이므로 $\xi_{\eq,j}(1-\xi_{\eq,j})=\dfrac{1}{1+e^{-z}}\cdot
\dfrac{e^{-z}}{1+e^{-z}}=\dfrac{e^{-z}}{(1+e^{-z})^2}$ 이고, 식~\eqref{eq:dxidz} 와 같다. 즉
$\dd\xi_{\eq,j}/\dd z=\xi_{\eq,j}(1-\xi_{\eq,j})$. 연쇄율 $\dd z/\dd V_n=s/w_j$ 를 곱하면
\begin{equation}
\frac{\dd\xi_{\eq,j}}{\dd V_n}=\frac{s\,\xi_{\eq,j}(1-\xi_{\eq,j})}{w_j}.
\label{eq:dxidV}
\end{equation}

이제 평형 도함수 식~\eqref{eq:dxidV} 를 관측식 식~\eqref{eq:dQdV} 의 peak 항에 대입한다.
\textbf{세 결과(각각 유도).} 식~\eqref{eq:dQdV} 의 peak 항 $Q_j\,\dd\xi_{\eq,j}/\dd V_n$ 으로부터:
\begin{itemize}[leftmargin=1.6em,itemsep=2pt]
\item \emph{위치}: $\dd\xi_{\eq,j}/\dd V_n\propto\xi_{\eq,j}(1-\xi_{\eq,j})$ 는 $\xi_{\eq,j}=\tfrac12$ 에서 최대다
  ($\xi(1-\xi)$ 가 $\xi=\tfrac12$ 서 최대). $\xi_{\eq,j}=\tfrac12\Leftrightarrow z=0\Leftrightarrow V_n=U_j$. 즉
  $V_\peak=U_j(T)$.
\item \emph{높이}: 그 점에서 $\xi_{\eq,j}(1-\xi_{\eq,j})=\tfrac14$ 이므로
  $\big(\dd Q/\dd V_n\big)_{\max}=C_\bg+Q_j/(4w_j)$($C_\bg$ 가 peak 폭서 거의 상수$\cdot$인접 전이 비중첩 가정), peak 순높이 $=Q_j/(4w_j)$.
\item \emph{너비}(FWHM): 아래 검산에서 $\Delta V_{1/2}=2w_j\ln(3+2\sqrt2)\approx3.53\,w_j$.
\end{itemize}
\begin{equation}
\boxed{\;V_\peak=U_j(T),\qquad \Big(\tfrac{\dd Q}{\dd V_n}\Big)^{\text{net}}_{\max}=\frac{Q_j}{4w_j},\qquad
\Delta V_{1/2}=2w_j\ln(3+2\sqrt2)\approx3.53\,w_j\;.}
\label{eq:eqpeak}
\end{equation}
\emph{면적 보존}: $\int(\dd Q/\dd V_n-C_\bg)\,\dd V_n=Q_j\int_0^1\dd\xi_{\eq,j}=Q_j$ 이므로 너비$\downarrow$ 면
높이$\uparrow$(높이$\times$너비 $\sim Q_j$). \emph{최종식과의 연결.} 이 평형 peak 의 세 양 $\{U_j,w_j,Q_j\}$ 가
최종 피팅식(\S\ref{sec:master})의 \emph{평형 상승부} 항을 이루며, 꼬리가 사라진 저율 OCV 에서 직접 닫힌다
(\S\ref{sec:synth} S1).

\begin{verifybox}
\textbf{FWHM 유도.} $g(z)=\xi_{\eq,j}(1-\xi_{\eq,j})$ 는 $z=0$ 서 최대 $\tfrac14$. 절반 $\tfrac18$ 되는 점:
$\xi(1-\xi)=\tfrac18$. $\xi=\tfrac12(1\pm\tfrac1{\sqrt2})$(이차방정식 $\xi^2-\xi+\tfrac18=0$ 의 해)이고,
$e^z=\xi/(1-\xi)=\dfrac{1+1/\sqrt2}{1-1/\sqrt2}$. 분자$\cdot$분모에 $\sqrt2$ 곱해 $\dfrac{\sqrt2+1}{\sqrt2-1}
=(\sqrt2+1)^2=3+2\sqrt2$. 따라서 $\Delta z=2\ln(3+2\sqrt2)$, $\Delta V_{1/2}=w_j\Delta z\approx3.53\,w_j$.
\end{verifybox}

\subsection{온도 의존(평형 기준선)}
\textbf{너비$\cdot$높이.} 이상 폭은 $w_j=RT/F$ 이므로 너비 $\propto T$ — \emph{저온서 좁아진다}($25^\circ$C 에서
$w=RT/F=8.314\times298/96485=25.7$ mV, $-15^\circ$C 에서 $8.314\times258/96485=22.2$ mV). 높이 $\propto1/w_j\propto
1/T$ — 저온서 높아진다.

\textbf{위치의 온도계수(Gibbs--Helmholtz 관계).} 전이 $j$ 의 \emph{삽입(리튬화) 반응} 자유에너지는 $\Delta G_j=-sF\,U_j$
다(평형 전위의 정의; 부호는 삽입 방향 기준). 온도로 미분하면 $\partial\Delta G_j/\partial T=-\Delta S_j$(Gibbs 관계)이므로
\begin{equation}
\frac{\partial U_j}{\partial T}=-\frac{1}{sF}\frac{\partial\Delta G_j}{\partial T}=\frac{\Delta S_j}{sF}.
\label{eq:dUdT}
\end{equation}
즉 peak \emph{위치}는 \emph{삽입 반응} 엔트로피 $\Delta S_j$ 에 따라 온도와 함께 이동한다(탈리튬화 기준으로 보면
$\Delta S_j\!\to\!-\Delta S_j$ 라 물리적으로 동치) \cite{reynier2004}(활성화 엔트로피 $\Delta S_{a,j}$ 와는 다른 양 —
후자는 \S\ref{sec:lag} 의 \emph{속도}에 들어간다).

\begin{keybox}
\textbf{평형 기준선의 한계.} 평형 peak 는 \emph{대칭}이고 저온서 \emph{더 좁고 높다}. 그러나 관찰은
``저온 $\to$ \emph{긴 꼬리}(넓고 비대칭)''로 \emph{정반대}이며, 평형 peak 는 전류(C-rate)와 무관한데 실제 peak
는 C-rate 에 민감하다. \textbf{따라서 꼬리의 주된 원인은 평형 열역학이 아니라, 전류가 흐를 때 평형을 못 따라가는
동역학이다.} 단, 그 동역학으로 가기 전에 — 평형 형태 자체가 \emph{이상 극한이 전부가 아니다}. 다음 절이 정규용액
상호작용($\Omega_j\ne0$)$\cdot$상분리까지 평형 기준선을 넓히고, 그 뒤 \S\ref{sec:lag} 가 동역학 지연을 다룬다.
\end{keybox}

% ====================================================================
\section{정규용액 상호작용과 상분리}\label{sec:regsol}
% ====================================================================
앞 절 평형 기준선은 \emph{이상 극한}($\Omega_j=0$)의 매끄러운 logistic 이었다. 그러나 흑연 staging 은 정규용액
상호작용 $\Omega_j\ne0$ 의 \emph{일차 상전이}라, 평형 형태 자체가 — 동역학에 앞서 — 비단조가 될 수 있다. 본 절은
그 상분리 열역학을 세운다: 그것이 (i) $\Omega_j>2RT$ 구간에서 peak 를 비단조로 만들고, (ii) 뒤에서 \emph{충방전
히스테리시스}(\S\ref{sec:hys_branch})의 분기를 만드는 뿌리다.

\subsection{정규용액 자유에너지 — 앞 절 $\mu$ 와 같은 $g_j$}
앞 절 화학퍼텐셜 식~\eqref{eq:mu} 의 상호작용 항 $\Omega_j(1-2\theta_j)$ 는 \emph{정규용액}(regular solution)
모형에서 온다. 전이 $j$ 의 진행률 $\xi$($=1-\theta_j$)에 대한 몰당 자유에너지는 혼합 entropy 와 상호작용 enthalpy
의 합이다:
\begin{equation}
g_j(\xi)=g_j^0+\underbrace{RT[\xi\ln\xi+(1-\xi)\ln(1-\xi)]}_{-T S_\mathrm{mix}}
+\underbrace{\Omega_j\,\xi(1-\xi)}_{\text{상호작용}}.
\label{eq:regsol_g}
\end{equation}
혼합 entropy 항은 앞 절 식~\eqref{eq:smix} 의 $S_\mathrm{mix}$ 그대로이고($\xi\leftrightarrow\theta$ 대칭), 상호작용
항은 \emph{평균장}에서 ``전환$\cdot$미전환'' 이웃 쌍의 분율 $\propto\xi(1-\xi)$ 에 쌍당 초과 에너지 $\Omega_j$
($\sim z[\epsilon_{AB}-\tfrac12(\epsilon_{AA}+\epsilon_{BB})]$, $z$=이웃 수)를 곱한 것이다 — $\Omega_j>0$ 이면
\emph{동종 자리 선호}(이종 쌍 불리)라 분리를 부추긴다. $g_j$ 의 점유율 미분이 곧 앞 절 식~\eqref{eq:mu} 의
$\mu_\mathrm{Li}(\theta)$ 다 — \emph{같은 $g_j$ 를 한 번은 peak(미분)로, 한 번은 상분리(곡률)로 본다}.

\subsection{곡률과 임계 온도 — 이중 우물}
상분리는 $g_j$ 의 \emph{곡률}이 정한다:
\begin{equation}
g_j''(\xi)=\frac{RT}{\xi(1-\xi)}-2\Omega_j.
\label{eq:regsol_gpp}
\end{equation}
entropy 항 $+RT/[\xi(1-\xi)]>0$ 은 항상 볼록(섞이려 함), 상호작용 항 $-2\Omega_j$ 는 $\Omega_j>0$ 라 항상 오목
(뭉치려 함). 양 끝($\xi\to0,1$)서는 entropy 가 발산해 $g_j''>0$ 이지만, 가운데서 $\Omega_j$ 가 크면 $-2\Omega_j$ 가
이겨 $g_j''<0$ — \emph{이중 우물}이다. 경계는 $g_j''$ 이 가장 작은 $\xi=\tfrac12$ 서 $4RT-2\Omega_j=0$, 곧
\textbf{임계 온도 $T_{c,j}=\Omega_j/2R$}: $T>T_{c,j}$($\Omega_j<2RT$)면 한 우물(고용체, 앞 절 logistic 그대로),
$T<T_{c,j}$($\Omega_j>2RT$)면 이중 우물(상분리).

\subsection{binodal $\cdot$ spinodal $\cdot$ 준안정}
이중 우물에서 한 균질 상보다 두 상($\xi$ 작은$\cdot$큰)으로 갈라지는 편이 자유에너지가 낮다 — 미스시빌리티 갭.
갈림은 두 경계로 규정된다:
\begin{itemize}[leftmargin=1.4em,itemsep=2pt]
\item \textbf{binodal} $\xi_{b,j}^\pm$ — 두 상이 \emph{평형}으로 공존하는 조성. 두 상의 화학퍼텐셜이 같고
  $\mu(\xi_b^-)=\mu(\xi_b^+)$ 교환이 자유에너지를 안 바꾸는 조건이 $g_j$ 의 \emph{공통 접선}(두 점서 같은 기울기
  $=\mu$$\cdot$같은 절편)으로 나타난다 — 공통 접선이 곧 평탄 전위(plateau)다.
\item \textbf{spinodal} $\xi_{s,j}^\pm$ — $g_j''=0$ 의 \emph{역학적 불안정} 경계. 안쪽($g_j''<0$)은 무한소 요동에도
  즉시 분리(불안정), binodal--spinodal 사이($g_j''>0$ 이나 공통 접선 위)는 \emph{준안정}.
\end{itemize}
조성축이 세 영역(안정 단일 상 / 준안정 / 불안정)으로 나뉘며 \emph{준안정}이 뒤 히스테리시스의 무대다. 준안정
영역에선 새 상의 \emph{핵}이 임계 크기를 넘어야(핵생성 장벽) 전이가 시작된다 — 장벽은 구동력이 작은 binodal
근처서 크고 안쪽으로 갈수록 준다. 장벽이 크면 계가 균질 준안정 분기를 따라 \emph{spinodal 까지 과주행}한 뒤에야
(거기선 준안정 최솟값 자체가 사라져 장벽 없이) 분리한다. \textbf{충$\cdot$방전이 서로 다른 끝에서 과주행하는 것이
충방전 히스테리시스의 근원}이다(\S\ref{sec:hys_branch}).

\textbf{전위로 보면 van der Waals 등온선.} 이중 우물이면 평형 전위 $V_{\eq,j}(\xi)$(뒤에서 $\mu_\mathrm{Li}/sF$ 로
정의)가 \emph{비단조} loop 를 그린다 — 압력--부피 van der Waals 등온선과 \emph{동형}이다($g_j\leftrightarrow$ Helmholtz
$A$, $\mu=\dd g/\dd\xi\leftrightarrow-P$, $\xi\leftrightarrow$ 부피). 음의 기울기 구간(spinodal 안쪽)이 vdW 의 양의
압축률($\partial P/\partial V>0$) 불안정 가지에, 공통 접선이 Maxwell 등면적 작도(=평탄 전위)에 대응한다.
\emph{이 loop 의 두 어깨가 곧 충$\cdot$방전 분기가 된다}(\S\ref{sec:hys_branch}) — 본 절 상분리가 뒤 히스테리시스로
이어지는 다리다.
\begin{flagbox}
\textbf{단일 입자 vs 다입자.} 위 spinodal 과주행은 \emph{단일 입자} 그림이다. 실제 전극은 입자 집단이라 입자마다
임계 핵생성이 조금씩 다른 전위에서 일어나 \emph{하나씩 차례로} 전환된다(Dreyer 다입자~\cite{dreyer2010}) —
같은 압력의 풍선 다발이 동시에 부풀지 않고 하나씩 부풀듯. 그래서 집단 거동은 단일 입자 과주행보다 분기 갈림을
\emph{매끄럽게$\cdot$좁게} 만들고, 단일 입자 spinodal 이 \emph{상한}, 실측은 그 아래(축소 인자 $\gamma_j$,
\S\ref{sec:hys_branch})에 놓인다.
\end{flagbox}

\subsection{peak 에의 영향, 그리고 다음}
$\Omega_j>2RT$ 면 isotherm 이 비단조라 그 구간 peak 는 앞 절 단일-phase logistic 로 기술되지 않고 \emph{상분리
plateau}(좁고 큰 peak)가 된다 — staging 의 날카로운 peak 한 원인. $\Omega_j\le2RT$ 면 spinodal 이 없어 앞 절
logistic 그대로다(상호작용은 폭만 보정). \textbf{이로써 평형 기준선이 이상($\Omega=0$)에서 정규용액($\Omega\ne0$)
까지 닫혔다.} 이제 전류가 흐를 때의 \emph{이탈}(동역학)로 간다.

% ====================================================================
\section{동역학 지연과 비대칭 꼬리}\label{sec:lag}
% ====================================================================
전류가 흐를 때 전환율 $\xi_j$ 는 평형 목표 $\xi_{\eq,j}$ 를 \emph{즉시} 따라가지 못한다. 이
\emph{지연}이 peak 를 어떻게 비대칭으로 늘리며, 왜 저온$\cdot$고율에서 꼬리가 길어지는가?

\subsection{완화 방정식의 유도}
\textbf{1차 완화의 근거.} 전이 $j$ 를 정/역 두 방향 속도 $r_j^+,r_j^-$ 를 갖는 2-상태 과정으로 보면, \emph{진행률}의
순 변화는 \emph{전향}(전이 완성=탈리튬화, $\xi_j$ 증가; $\propto$ 아직 전이 안 된 분율 $1-\xi_j=\theta_j$, 즉 잔여 Li
점유율)과 \emph{역향}(재리튬화, $\xi_j$ 감소; $\propto$ 이미 전이된(빈자리) 분율 $\xi_j$)의 차다:
\begin{align}
\frac{\dd\xi_j}{\dd t}&=r_j^+(1-\xi_j)-r_j^-\xi_j
=r_j^+-(r_j^++r_j^-)\xi_j\notag\\
&=(r_j^++r_j^-)\Big[\frac{r_j^+}{r_j^++r_j^-}-\xi_j\Big]
\equiv k_j\,(\xi_{\sstat,j}-\xi_j),
\label{eq:relax}
\end{align}
$k_j\equiv r_j^++r_j^-$(완화 속도), $\xi_{\sstat,j}\equiv r_j^+/(r_j^++r_j^-)$(정상점). 유일한 가정은 ``$r_j^\pm$ 를
$\xi_j$ 의 작은 변화에 대해 국소 상수로 본다''이며, 그 위에서 위 재배열은 \emph{근사가 아니라 대수적 항등}이다.
\emph{단, 전이 전체를 \textbf{단일 완화 모드}(1차 ODE)로 보는 것 자체가 ansatz 다} — 핵생성$\cdot$상경계 이동 등
여러 시간척도를 하나의 유효 모드로 coarse-grain 한 것으로, 한 완화 시간이 지배할 때 타당하다(다중 모드$\cdot$입자
분산 효과는 \S\ref{sec:dist}).
\textbf{정상점과 평형값의 일치}($\xi_{\sstat,j}=\xi_{\eq,j}$)는 detailed balance 의 귀결인데, 그 검증
($r_j^+/r_j^-=e^{\mathcal A_j/RT}$ 가 식~\eqref{eq:logistic} 을 재생함)은 \S\ref{sec:barrier} 에서 정/역 속도를
구체화할 때 이루어진다 — 여기서는 $\xi_{\sstat,j}=\xi_{\eq,j}$ 로 둔다(전방 참조).

\textbf{속도상수와 활성화 배리어 (Eyring).} 전이상태 이론에서 속도상수는, 전이상태를 통과하는 보편 빈도
$k_BT/h$ 에 장벽을 넘을 Boltzmann 인자를 곱한 것이다 \cite{eyring1935}:
\begin{equation}
k_j\simeq k_0\,\exp\!\Big[-\frac{\Delta G_{a,j}}{RT}\Big]\ \ (k_0=k_BT/h,\ \text{현상론적 prefactor}),\qquad \Delta G_{a,j}=\Delta H_{a,j}-T\Delta S_{a,j}.
\label{eq:eyring}
\end{equation}
여기 $k_j$ 는 \emph{완화 속도}($=r_j^++r_j^-$, 식~\eqref{eq:relax}) \emph{그대로}이며, 그 \emph{척도}가 활성화 배리어의
Arrhenius/Eyring 형임을 적은 것이다(그래서 ``$=$''가 아니라 ``$\simeq$''이고, $k_0$ 는 $O(1)$ 수 인자$\cdot$교환전류$\cdot$
활동도를 흡수하는 현상론 prefactor — \S\ref{sec:barrier}). 정/역 분해는 \S\ref{sec:barrier} 의 $r_j^\pm$ 로 구체화되며,
forward 우세 꼬리에서 $k_j\simeq r_j^+$ 가 된다(peak 중심부 $\mathcal A_j\approx0$ 에서는 $r_j^+\approx r_j^-$ 라 합이 인자
$2$ 만큼 크지만 그 수 인자도 $k_0$ 에 흡수). \emph{$k_j$ 를 forward 속도로 \textbf{재정의하지 않는다}.}
\textbf{저온서 $k_j$ 가 작아진다}($e^{-\Delta H_{a,j}/RT}$ 가 급감) — 이것이 곧 ``저온 긴 꼬리''의 근원이다.
(꼬리 길이 \emph{절대값}은 $L_{q,j}\propto1/k_0$ 라 $k_0$ 에 의존하나, 본 장 \emph{결론}(꼬리의 온도$\cdot$전위
\emph{의존성}과 부호, Arrhenius 기울기로 뽑는 $\Delta H_a$)은 $k_0$ 를 피팅 prefactor 로 흡수하므로 그 값에
무관하다 — 단 흡수 인자들이 피팅 온도창에서 $T$-무관이라는 전제 하이며, 강한 $T$ 의존은 Arrhenius
절편$\cdot$기울기에 계통오차로 들어간다.)

\subsection{지연이 만드는 비대칭 꼬리}
\textbf{측정축으로의 변환(시간 $\to$ 용량).} 정전류 방전에서 무차원 용량 좌표 $q=\int|I|\,\dd t/Q_\cell$ 는
$\dd q/\dd t=|I|/Q_\cell$(단위시간당 흘린 용량분율)다. 연쇄율 $\dd\xi_j/\dd q=(\dd\xi_j/\dd t)/(\dd q/\dd t)$ 로
식~\eqref{eq:relax} 을 $q$ 미분으로 바꾸면
\begin{equation}
\frac{\dd\xi_j}{\dd q}=\frac{k_j(\xi_{\eq,j}-\xi_j)}{|I|/Q_\cell}=\frac{Q_\cell\,k_j}{|I|}\,(\xi_{\eq,j}-\xi_j).
\label{eq:dxidq}
\end{equation}

\textbf{지연의 방정식과 그 해.} peak 를 지난(post-peak) 영역에서는 목표가 logistic 양극단으로 포화해
$\xi_{\eq,j}(1-\xi_{\eq,j})\to0$, 따라서 $\dd\xi_{\eq,j}/\dd q\simeq0$ 이다. 지연을 $r_j\equiv\xi_{\eq,j}-\xi_j$
(첨자 없는 $r_j$=지연; 정/역 \emph{속도} $r_j^\pm$ 와 구별)라 하면 $\dd r_j/\dd q=\dd\xi_{\eq,j}/\dd q-\dd\xi_j/\dd q\simeq-\dd\xi_j/\dd q$ 이고, 식~\eqref{eq:dxidq} 을
넣으면
\begin{equation}
\frac{\dd r_j}{\dd q}=-\frac{Q_\cell k_j}{|I|}\,r_j\equiv-\frac{1}{L_{q,j}}\,r_j,\qquad
\boxed{\;L_{q,j}=\frac{|I|}{Q_\cell\,k_j}\;}.
\label{eq:Lq}
\end{equation}
이 1계 선형 ODE 의 해는 $q_a$(꼬리 시작점 — 조작적으로 평형 항이 포화해 $\dd\xi_{\eq,j}/\dd q$ 가 정점값의 일정
분율(예 $\sim5\%$) 이하로 떨어지는 \emph{단일} 좌표 = 상승부에서 꼬리로 넘어가는 곳; 정확한 창은 피팅에서 선택, \S\ref{sec:synth} S1)에서 적분해
\begin{equation}
r_j(q)=r_j(q_a)\,\exp\!\Big[-\frac{q-q_a}{L_{q,j}}\Big],
\label{eq:rsol}
\end{equation}
즉 지연이 지수로 감쇠한다. \textbf{$L_{q,j}$ 는 ``지연이 $1/e$ 로 줄어드는 방전 거리''}, 곧 꼬리의 특성 길이다.
관측 peak 항은 $\dd\xi_j/\dd q=-\dd r_j/\dd q=(r_j(q_a)/L_{q,j})\exp[-(q-q_a)/L_{q,j}]$ 로 같은 지수형이다.

\textbf{용량축 $\to$ 전압축.} post-peak 에서 $V_n$ 이 $q$ 에 거의 선형이면($\dd V_n/\dd q$ 가 $q_a$ 둘레 상수;
$V_a\equiv V_n(q_a)$), $q-q_a\simeq(V_n-V_a)/(\dd V_n/\dd q)$ 이고, 꼬리 진행 방향에서 $(V_n-V_a)$ 와 $\dd V_n/\dd q$
는 \emph{같은 부호}라 그 비 $=|V_n-V_a|/|\dd V_n/\dd q|\ge0$ 이다. (staging plateau 처럼 $\dd V_n/\dd q\to0$ 인 2-상
공존 구간은 $L_{V,j}\to0$ 이라 단일-phase 가정 밖이다 — 비고$\cdot$\S\ref{sec:falsify}; 본 꼬리식은 plateau 를
벗어난 단상 영역에 쓴다.) 따라서 꼬리는 전압축에서 감쇠한다(\emph{단방향}
— post-peak 전향 영역 $s(V_n-V_a)\ge0$ 에서만 정의되며, 양방향 대칭 감쇠가 아니다):
\begin{equation}
\frac{\dd Q}{\dd V_n}\Big|_{\text{tail}}\propto\exp\!\Big[-\frac{|V_n-V_a|}{L_{V,j}}\Big]\quad\big(s(V_n-V_a)\ge0\big),\qquad
L_{V,j}=\Big|\frac{\dd V_n}{\dd q}\Big|_{q_a}L_{q,j}.
\label{eq:tail}
\end{equation}

\begin{verifybox}
\textbf{$L_{q,j}$ 무차원.} $|I|$=$\mathrm{C/s}$, $Q_\cell$=$\mathrm C$, $k_j$=$\mathrm{s^{-1}}$ 이므로
$|I|/(Q_\cell k_j)=(\mathrm{C/s})/(\mathrm C\cdot\mathrm{s^{-1}})=1$ — $q$ 와 같은 무차원 길이. \textbf{$L_{q,j}$ 상수는
\emph{국소 근사}다}: 정확한 지연 해는
\[ r_j(q)=r_j(q_a)\exp\!\Big[-\int_{q_a}^q\frac{\dd q'}{L_{q,j}(q')}\Big] \]
이고, 꼬리 감쇠 길이 안에서 구동력 $\mathcal A_j$(따라서 $k_j$)가 천천히 변해 $|\dd L_{q,j}/\dd q|\ll1$ 일 때만 $L_{q,j}$ 를
그 구간 상수로 보고 적분 밖으로 뺀다(slowly-varying). 급변 구간에서는 단일 지수가 깨져 \S\ref{sec:dist} 의 중첩으로 넘어간다.
\end{verifybox}

\subsection{온도$\cdot$C-rate 의존 — 관찰과의 일치}
식~\eqref{eq:Lq},~\eqref{eq:eyring} 에서 $1/k_j\simeq(h/k_BT)e^{\Delta G_{a,j}/RT}$ 이므로($k_0=k_BT/h$ 현상론)
\begin{equation}
L_{q,j}\simeq\frac{|I|}{Q_\cell}\frac{h}{k_BT}\exp\!\Big[\frac{\Delta G_{a,j}}{RT}\Big]
\;\Longrightarrow\;
\begin{cases}
\text{저온}\ (T\downarrow): & k_j\downarrow\Rightarrow L_{q,j}\uparrow\Rightarrow\text{꼬리 길어짐}\ (\checkmark),\\[2pt]
\text{고율}\ (|I|\uparrow): & L_{q,j}\uparrow\Rightarrow\text{꼬리 길어짐, peak 낮아짐}.
\end{cases}
\label{eq:tailTC}
\end{equation}

\textbf{단일지수에서 늘어진(stretched) 꼬리로 — 후반부의 동기.} 위 꼬리는 \emph{한 입자$\cdot$한 배리어}
$\Delta G_{a,j}$ 가정의 \emph{단일 지수}다. 그러나 실제 전극의 꼬리는 그보다 \emph{늘어져}(stretched) 관측된다 —
입자마다 크기$\cdot$결정성$\cdot$국소 staging 이 달라 활성화 배리어가 \emph{분포}를 이루고, 서로 다른 감쇠 길이
$L_{q,j}$ 가 겹치기 때문이다. 이 늘어짐을 정량화하는 것이 후반부(\S\ref{sec:stattools}$\cdot$\S\ref{sec:dist})의
과제이며, 여기서 얻은 지수 꼬리 항은 최종식(\S\ref{sec:master})의 \emph{동역학 항}이 된다(다율$\cdot$다온도로
닫힌다 — \S\ref{sec:synth} S3$\cdot$S4).

\begin{keybox}
\textbf{동역학 지연의 요지.} 동역학 지연은 평형 peak 의 \emph{뒤쪽}에 지수 꼬리를 붙여 peak 를 \emph{비대칭}으로 만들고, 면적
보존상 \emph{높이를 낮춘다}. (수치 감: $\Delta H_{a,j}\approx40$ kJ/mol 이면 $25^\circ$C$\to-15^\circ$C 에서 꼬리
길이 $L_{q,j}$ 가 약 $10$배 이상 길어진다 — $e^{(\Delta H_a/R)(1/258-1/298)}\approx e^{2.5}\approx12$ 에 prefactor
$1/T$ 까지.) 빠른 kinetics 극한($k_j\to\infty$)에서는 $L_{q,j}\to0$ 이라 꼬리가 사라지고 평형
peak(\S\ref{sec:eqpeak})로 매끄럽게 복귀한다. \emph{온도와 C-rate 의존성은 위 식들로 설명됐다; 남은 것은 전위가
$k_j$ 를 바꾸는 방식이다 — 다음 절.}
\end{keybox}

% ====================================================================
\section{전위에 의한 유효 배리어}\label{sec:barrier}
% ====================================================================
관찰의 핵심은 ``온도에 의한 배리어 \emph{외에}, 극판 전위가 배리어를 추가로 낮춘다''는 것이었다.
전위는 식~\eqref{eq:eyring} 의 $\Delta G_{a,j}$ 에 어떻게 들어가는가?

\subsection{구동력}
\S\ref{sec:stage} 에서 구동 전위 $V_\drive$(기본 $=V_n$)를 두었다. 전이 $j$ 의 \textbf{구동력}(affinity)은
구동 전위가 전이 중심에서 벗어난 정도로, 식~\eqref{eq:eqcond} 의 평형 구동과 같은 형태의 몰당 자유에너지다(기본형
$V_\drive=V_n$ 에서 $\mathcal A_j/RT=s(V_n-U_j)/w_j$ — 식~\eqref{eq:logistic} 의 logit 구동항과 동일):
\begin{equation}
\mathcal A_j=sF\,(V_\drive-U_j).
\label{eq:affinity}
\end{equation}

\subsection{유효 배리어 — Butler--Volmer 방향성 장벽 이동}
\textbf{구동력에 의한 forward 장벽 강하.} 정/역 두 방향(\S\ref{sec:lag} 의 $r_j^\pm$)을 보면, 구동력 $\mathcal A_j$ 는
두 방향의 활성화 장벽을 \emph{비대칭}으로 옮긴다 — 표준 Butler--Volmer/Brønsted--Evans--Polanyi 선형 자유에너지
관계다 \cite{bardfaulkner,evanspolanyi1938}. 전달 계수 $\chi\in[0,1]$(구동력 중 forward 장벽을 낮추는 몫; 반응
좌표상 전이상태의 위치로 해석)로 — 정$\cdot$역 반응은 \emph{같은 전이상태}를 공유하므로 전이상태 통과 빈도
prefactor $k_0=k_BT/h$(식~\eqref{eq:eyring})가 두 방향에 \emph{공통}이다 —
\begin{equation}
r_j^+=k_0\,e^{-(\Delta G_{a,j}-\chi\mathcal A_j)/RT},\qquad
r_j^-=k_0\,e^{-(\Delta G_{a,j}+(1-\chi)\mathcal A_j)/RT},\qquad k_0=\frac{k_BT}{h}.
\label{eq:bv}
\end{equation}

\textbf{전달계수 분할과 합 1의 근거.} 전달 계수 $\chi$ 는 전이상태의 \emph{반응 좌표상 분율
위치}다($0$=반응물 쪽, $1$=생성물 쪽, $\tfrac12$=가운데). 구동력 $\mathcal A_j$ 가 두 우물의 상대 높이를
기울이면, \emph{같은 한 점}인 전이상태가 그 분율만큼 따라 움직여 — forward 장벽(반응물$\to$전이상태)은
$\chi\mathcal A_j$ 만큼 내려가고, backward 장벽(생성물$\to$전이상태)은 나머지 $(1-\chi)\mathcal A_j$ 만큼 올라간다
\cite{evanspolanyi1938,bardfaulkner}. \textbf{두 계수의 합이 정확히 $1$ 인 것은 자유 선택이 아니라 detailed
balance 의 강제}다: 아래에서 정/역 비 $r_j^+/r_j^-=e^{[\chi+(1-\chi)]\mathcal A_j/RT}$ 가 자유에너지 차의
Boltzmann 인자 $e^{\mathcal A_j/RT}$ 와 같으려면 지수의 두 몫의 합이 \emph{반드시} $1$ 이어야 한다(아니면
운동학이 주는 평형 상수가 열역학 구동력과 어긋나 평형이 깨진다). 즉 $\chi$ 자체는 $0\!\sim\!1$ 의 자유
파라미터지만, ``$\chi$ 와 $1-\chi$ 로 갈라져 합이 $1$''은 두 방향이 \emph{같은 전이상태를 공유}함의 대수적
귀결이다.

\textbf{평형 불변 (detailed balance).} 위 합$=1$ 로 두 방향 비는
\begin{equation}
\frac{r_j^+}{r_j^-}=\exp\!\Big[\frac{(\chi\mathcal A_j)-(-(1-\chi)\mathcal A_j)}{RT}\Big]
=\exp\!\Big[\frac{\mathcal A_j}{RT}\Big],
\label{eq:db}
\end{equation}
($\chi$ 가 상쇄). 따라서 정상점은 $\xi_{\sstat,j}=\dfrac{r_j^+}{r_j^++r_j^-}=\dfrac{1}{1+r_j^-/r_j^+}
=\dfrac{1}{1+e^{-\mathcal A_j/RT}}$. 기본형 $V_\drive=V_n$ 에서 $\mathcal A_j=sF(V_n-U_j)$, $w_j=RT/F$ 이므로 이는
\begin{equation}
\xi_{\sstat,j}=\frac{1}{1+\exp[-s(V_n-U_j)/w_j]}=\xi_{\eq,j}
\label{eq:sseq}
\end{equation}
즉 \S\ref{sec:eqpeak} 의 logistic~\eqref{eq:logistic} 과 \emph{정확히 일치}한다 — \S\ref{sec:lag} 에서 전방참조로
둔 $\xi_{\sstat,j}=\xi_{\eq,j}$ 가 여기서 \emph{증명}된다. \textbf{장벽 이동은 평형 target 을 옮기지 않는다}: 평형은 두
상태의 \emph{자유에너지 차}($\mathcal A_j$)가 정하고, 장벽(활성화)은 \emph{거기 도달하는 속도}만 정한다(전이상태
이론의 표준 역할 분리; $\chi$ 가 식~\eqref{eq:db} 에서 상쇄되는 것이 그 표현).

\textbf{완화 속도의 유효 배리어(regime 명시).} 완화 속도는 $k_j=r_j^++r_j^-$ 다. 방전(탈리튬화)에서 전이가
\emph{완성}되는 꼬리 영역은 구동 전위가 전이 중심을 forward 방향으로 지난 곳이므로($\xi_j$ 가 1 로 갈 때
$V_n>U_j$, 따라서 $\mathcal A_j=sF(V_n-U_j)>0$, $s=+1$), 식~\eqref{eq:db} 로 $r_j^+/r_j^-=e^{\mathcal A_j/RT}$ 다.
이 비가 충분히 크면($\mathcal A_j\gtrsim RT$, forward 우세) $r_j^-$ 를 무시해 $k_j\simeq r_j^+$ — 상대 보정은
$r_j^-/r_j^+=e^{-\mathcal A_j/RT}$ 라 $\mathcal A_j=RT$ 서 약 $37\%$($3RT$ 서 $\sim5\%$)이므로, 단일지수 꼬리는 $\mathcal A_j$ 가
\emph{수 $RT$($\gtrsim3RT$, 보정 $\lesssim5\%$)} 이상일 때 정확하다:
\begin{equation}
\boxed{\;\Delta G_{\eff,j}=\Delta G_{a,j}-\chi\,\mathcal A_j\;,\qquad
k_j\simeq k_0\exp\!\Big[-\frac{\Delta G_{a,j}-\chi\mathcal A_j}{RT}\Big]\quad(\mathcal A_j\gtrsim 3RT)\;.}
\label{eq:keff}
\end{equation}
구동력이 클수록($\mathcal A_j>0$) forward 유효 배리어가 낮아져 완화 속도 $k_j$ 가 커진다 — \emph{평형 target 은
\S\ref{sec:eqpeak} 그대로}. \textbf{이것이 ``전위가 배리어를 낮춘다''의 정확한 자리다(평형이 아니라 속도).}
\emph{유효 작동 창}: $\mathcal A_j\approx0$ 인 peak \emph{중심부}에서는 $r_j^+\approx r_j^-$ 라 $k_j\simeq2r_j^+$ 로
단일지수 형태가 깨지지만, 거기서는 관측이 평형 상승부(\S\ref{sec:eqpeak})로 지배되므로, 단일지수 \emph{꼬리}
근사는 forward 우세 확장 꼬리($\mathcal A_j\gtrsim RT$)에 쓴다. (파라미터 식별: $\{\Delta G_{a,j},\chi,k_0\}$ 피팅$\cdot$
$V_\drive=V_n$ 측정$\cdot$$U_j$ 열역학 입력 — \S\ref{sec:synth} S1–S4.)
\begin{boundbox}
\textbf{선형의 한계(Marcus 유추)와 작동 창.} $-\chi\mathcal A_j$ 는 \emph{소구동력 1차} 근사다. Marcus 포물면
$\Delta G^\ddagger(\mathcal A)=\dfrac{(\lambda-\mathcal A)^2}{4\lambda}=\dfrac{\lambda}{4}-\dfrac{\mathcal A}{2}
+\dfrac{\mathcal A^2}{4\lambda}$ 의 $\mathcal A=0$ 1차 전개는 계수 $\partial\Delta G^\ddagger/\partial\mathcal A|_0
=-\tfrac12$, 즉 대칭 극한 $\chi=\tfrac12$ 을 \emph{유추}로 준다 \cite{marcus1956}. $|\mathcal A_j|\sim\lambda$(큰
과전위$\cdot$역전 영역)에서는 2차 곡률로 깨진다. \emph{intercalation 상전이로의 1차 근거는 BV/BEP 경험식이고,
Marcus 는 그 선형 형태의 물리적 유추로만 인용한다.} \textbf{본 장의 유효 작동 창은 $RT\lesssim\mathcal A_j\ll\lambda$}
— forward 우세로 $k_j\simeq r_j^+$ 가 서고(아래 한계), 선형 배리어 낮춤이 Marcus 2차로 깨지기 전이며,
$\lambda\gg RT$ 라 이 창이 존재한다.
\end{boundbox}
\begin{boundbox}
\textbf{속도 제한(직렬 시간척도; 불연속 절단을 피하는 구현).} 전하전달이 빨라도 수송$\cdot$유한 자리$\cdot$
고체 확산이 전체 진행을 제한할 수 있다. 임의 $\min$/clip 대신 율속 단계의 직렬 합 $1/k_j^{\eff}=1/k_{j,\mathrm{act}}+1/k_{j,\lim}$
으로 둔다($k_{j,\mathrm{act}}$=식~\eqref{eq:keff} 의 활성화 속도; $k_{j,\lim}\to\infty$ 면 $k_j^{\eff}\to k_{j,\mathrm{act}}$ 로 활성화
지배 환원). 이로써 강구동서 $\Delta G_{\eff,j}<0$ 가 되어도 $\max(\cdot,0)$ 같은 절단을 쓰지 않는다. 이하 꼬리식($L_{q,j}$)의
속도는 이 $k_j^{\eff}$ 다(활성화 지배면 $\simeq k_{j,\mathrm{act}}$; $k_{j,\lim}$ 의 $V_\drive$/SOC 의존 가능성은 \S\ref{sec:falsify}
(i) 수송 진단으로 점검). (이 $\eff$ 첨자는 \emph{직렬 율속}
합성이고, 식~\eqref{eq:keff} 의 $\Delta G_{\eff,j}$ 의 $\eff$ 는 \emph{전위에 의한 배리어 낮춤}으로 — 같은 첨자지만 다른
물리다.)
\end{boundbox}

\subsection{전위$\cdot$C-rate 가 peak 를 바꾸는 통로}
식~\eqref{eq:keff} 를 식~\eqref{eq:Lq} 에 넣으면 $L_{q,j}\propto1/k_j^{\eff}$ 라 $\partial\ln L_{q,j}/\partial V_\drive=-\partial\ln k_j^{\eff}/\partial V_\drive$ 이고($1/k_j^{\eff}=1/k_{j,\mathrm{act}}+1/k_{j,\lim}$,
$k_{j,\lim}$ 은 $V_\drive$ 무관 가정), $\partial\mathcal A_j/\partial V_\drive=sF$(식~\eqref{eq:affinity})이고 전위 의존은
$k_{j,\mathrm{act}}$ 만 통하므로(연쇄율 $\partial\ln k_j^{\eff}/\partial V_\drive=\tfrac{k_{j,\lim}}{k_{j,\mathrm{act}}+k_{j,\lim}}\,\partial\ln k_{j,\mathrm{act}}/\partial V_\drive$)
꼬리 길이의 \emph{전위 의존}은
\begin{equation}
\frac{\partial\ln L_{q,j}}{\partial V_\drive}=-\frac{k_{j,\lim}}{k_{j,\mathrm{act}}+k_{j,\lim}}\,\frac{\chi sF}{RT}\le0\quad(s=+1).
\label{eq:LqV}
\end{equation}
활성화 지배 $k_{j,\mathrm{act}}\ll k_{j,\lim}$ 극한에서 인자 $k_{j,\lim}/(k_{j,\mathrm{act}}+k_{j,\lim})\to1$ 이라 $-\chi sF/RT$ 로
환원된다; 유한 공율속이면 그만큼 \emph{감쇠}하나 인자 $\in(0,1]$ 이라 \emph{부호}($\chi>0$ 이면 $<0$ — 전위 의존)는
그대로다(\S\ref{sec:falsify} 반증 지문이 율속 비율에 견고한 까닭).
\begin{keybox}
\textbf{세 인자의 peak 진입.} 구동 전위가 커지면 forward 유효 배리어가 낮아져 꼬리가 \emph{짧아진다}. 또 측정축에서는 분극
$s_I|I|R_n$(식~\eqref{eq:vapp})이 apparent peak 위치를 밀어, C-rate 가 \emph{위치} 이동에도 기여한다. \emph{이로써
세 인자(온도$\cdot$C-rate$\cdot$전위)가 모두 $k_j$ 와 $L_{q,j}$ 를 통해 peak 에 들어왔다.} 남은 물리는 하나 —
전극 입자가 \emph{동일하지 않다}는 것이다. 그 비균질이 \S\ref{sec:lag} 에서 예고한 \emph{늘어진}(stretched) 꼬리를
만든다. 이를 다루려면 분포$\cdot$가중 평균$\cdot$분산 전파의 통계 도구가 필요하므로 — \emph{다음 절
(\S\ref{sec:stattools})은 잠시 그 도구를 정비}하고, 그 위에서 \S\ref{sec:dist} 가 입자 분포를 정량화한다.
\end{keybox}

% ====================================================================
\section{입자 분포를 다루는 통계적 기술}\label{sec:stattools}
% ====================================================================
\textbf{도구 정비 막간.} 앞 절들이 \emph{한 입자}의 물리를 닫았다면, 여기서 잠시 멈춰 \emph{입자 집단}을 다룰
통계 도구를 벼린다 — \S\ref{sec:lag} 가 예고하고 \S\ref{sec:barrier} 가 가리킨 \emph{늘어진 꼬리}가 이 도구로
정량화된다.
\S\ref{sec:barrier} 까지 한 전이를 활성화 배리어 \emph{한 값} $\Delta G_{a,j}$ 로 다뤘으나, 실제 전극은
입자 집단마다 배리어가 다르다(\S\ref{sec:dist}). 따라서 관측은 \emph{배리어 분포}에 대한 합이며, 이를 기술하려면
분포$\cdot$가중 평균$\cdot$분산 전파$\cdot$변수 변환$\cdot$연속 중첩의 다섯 통계적 기술이 필요하다. 본 절은 이를
확률과 미적분만으로 본 문제에 맞게 \emph{자족적으로} 확립한다 — 이후 \S\ref{sec:dist} 가 그대로 인용하며, 본 절의
결과식은 \S\ref{sec:dist} 의 \emph{입력}이다.

\subsection{활성화 배리어의 분포와 밀도}
전극 입자 집단은 크기$\cdot$결정성$\cdot$국소 staging 이 제각각이라 활성화 배리어 $G$ 가 단일 값이 아니라 분포를
이룬다. 이를 밀도(density) $\rho_G(G)$ 로 기술하며, $\rho_G(G)\,\dd G$ 는 배리어가 구간 $[G,\,G+\dd G]$ 에 드는
집단의 분율이다:
\begin{equation}
\rho_G(G)\,\dd G=\big(\text{배리어가 }[G,\,G+\dd G]\text{ 에 드는 집단의 분율}\big),\qquad [\rho_G]=1/[\text{에너지}].
\label{eq:density}
\end{equation}
전 구간의 분율 합이 1 이므로 정규화(normalization) 조건은
\begin{equation}
\int_0^\infty\rho_G(G)\,\dd G=1.
\label{eq:norm}
\end{equation}
본 장의 $\rho_G$ 는 단순 입자 수 밀도가 아니라, 각 집단이 관측 $\dd Q$ 에 기여하는 용량 몫까지 반영한 \emph{용량 가중 밀도}(capacity-weighted density)다.
따라서 $\rho_G$ 가중 합이 곧 관측 신호가 된다.

\subsection{분포에 대한 가중 평균(ensemble average)}
관측량은 한 집단이 아니라 모든 집단의 기여를 분율로 가중해 합한 값이다. 집단 양 $f(G)$ 의 가중 평균은 합을
적분으로 옮겨
\begin{equation}
\langle f\rangle=\int_0^\infty f(G)\,\rho_G(G)\,\dd G,
\label{eq:wavg}
\end{equation}
즉 각 집단 값 $f(G)$ 를 분율 $\rho_G\,\dd G$ 로 가중한 적분이다. \S\ref{sec:dist} 의 중첩식(식~\eqref{eq:superpose})
은 각 배리어 집단의 꼬리 기여를 $f(G)$ 로 두고 이 평균을 취한 것으로, 임의의 결합이 아니라 분율 가중 합이다.

\subsection{분산의 전파와 지수적 증폭}
분포는 평균뿐 아니라 그 폭으로도 특징지어진다. 분산(variance) $\sigma_G^2=\langle(G-\langle G\rangle)^2\rangle$,
표준편차(standard deviation) $\sigma_G$ 가 그 척도이며, 선형 척도 변환의 분산은 $\mathrm{Var}(aG)=a^2\sigma_G^2$ 다.
꼬리 길이는 배리어의 지수이므로 — 동일 $T\cdot|I|\cdot$구동력에서 집단 간 차이가 배리어 $G$ 뿐이라
$L_q\propto e^{(G-\chi\mathcal A_j)/RT}$(\S\ref{sec:barrier}$\cdot$\S\ref{sec:dist}; 구동항 $\chi\mathcal A_j$ 는 집단
\emph{공통}이라 아래 const 에 흡수) — 로그를 취하면 $\ln L_q=\text{const}+G/RT$ 로 선형이 되어($\chi\mathcal A_j/RT$
는 그 const 에 들어가 \emph{분산에 무관}), $a=1/RT$ 로
\begin{equation}
\mathrm{Var}(\ln L_q)=\Big(\tfrac{1}{RT}\Big)^2\sigma_G^2,\qquad \sigma_{\ln L_q}=\frac{\sigma_G}{RT}.
\label{eq:varprop}
\end{equation}
배리어 분산이 길이(로그) 분산으로 $1/RT$ 만큼 증폭되므로, 저온($RT\downarrow$)일수록 길이 분산이 커진다. 고정된
입자 분포가 온도를 통해 거동을 바꾸는 메커니즘이며, \S\ref{sec:dist} 의 ``저온에서 꼬리가 더 늘어진다''의 근원이다.

\subsection{변수 변환과 밀도 보존}
분포의 가로축을 배리어 $G$ 에서 길이 $L$ 로 옮길 때 한 집단의 분율은 좌표 선택에 무관하게 보존되므로, 두 밀도는
\begin{equation}
A_L(L)\,\dd L=\rho_G(G)\,\dd G\quad\Longrightarrow\quad A_L(L)=\rho_G\big(G(L)\big)\,\Big|\frac{\dd G}{\dd L}\Big|
\label{eq:jacobian}
\end{equation}
로 연결된다($|\dd G/\dd L|$=Jacobian). 본 문제에 대한 구체적 전개는 \S\ref{sec:dist} 의 심화에서 수행하며,
본문에는 분율 보존만으로 충분하다.

\subsection{지수 감쇠의 연속 중첩}
빠른 집단(작은 $L$)은 peak 근처에서 감쇠를 마치고 느린 집단(큰 $L$)은 꼬리를 멀리까지 연장한다. 단일 지수 $e^{-x/L}$ 를 여러 $L$
에 대해 분율로 가중 중첩하면(식~\eqref{eq:wavg}), 결과는 단일 지수가 아니라 초기에 빠르고 후미에서 완만한
stretched exponential 꼬리가 된다 \cite{lindsey1980,johnston2006,svare2000}.

\begin{keybox}
\textbf{요약.} 입자 분포는 밀도 $\rho_G$(정규화 식~\eqref{eq:norm})로, 관측은 그에 대한 가중 평균
$\int f\rho_G\,\dd G$(식~\eqref{eq:wavg})로 기술된다. 길이가 배리어의 지수이므로 배리어 분산은 길이 분산으로
$\sigma_G/RT$ 만큼 증폭되고(식~\eqref{eq:varprop}) 저온에서 꼬리가 더 늘어진다. 배리어 분포는 길이 분포로 분율
보존 변환되며(식~\eqref{eq:jacobian}), 지수 감쇠의 연속 중첩은 stretched exponential 꼬리를 준다. \S\ref{sec:dist}
는 이 다섯 기술로 입자 분포를 정량화한다 — 새 물리가 아니라 분율 가중 합산이다.
\end{keybox}

% ====================================================================
\section{배리어 분포와 중첩}\label{sec:dist}
% ====================================================================
지금까지 한 전이의 배리어 $\Delta G_{a,j}$ 를 \emph{한 값}으로 다뤘다. 그러나 실제 전극의 입자는
크기$\cdot$국소 staging$\cdot$결정성$\cdot$결정립계가 제각각이다. 이 차이를 무시할 수 없다 — 관측 peak/꼬리는
\emph{여러 입자 집단의 합}이기 때문이다. 어떻게 합치는가?

\textbf{입자 분포 $\to$ 배리어 분포 $\to$ 중첩(네 단계).}
\begin{enumerate}[label=\textbf{(\alph*)},leftmargin=2.2em,itemsep=2pt]
\item \textbf{입자별 배리어 차이.} 큰 입자/작은 입자, 결정성이 좋은/나쁜 입자, staging 환경이 다른 도메인은
  같은 전이라도 넘어야 할 활성화 장벽이 다르다.
\item \textbf{배리어 분포.} 한 값 대신 \emph{분포} $\rho_G(G)$ 로 기술한다(\S\ref{sec:stattools} 의 밀도;
  정규화 식~\eqref{eq:norm} $\int_0^\infty\rho_G(G)\,\dd G=1$). $\rho_G$ 는 입자 크기 분포$\cdot$국소 상태에서
  오는 \emph{독립} 물리량이며, 각 집단이 관측 $\dd Q$ 에 내는 \emph{용량 몫}까지 흡수한 \emph{용량 가중} 분포다. (구동력 $\mathcal A_j$ 는 집단 공통 —
  동일 평균장 $V_n$ — 이라 분포는 배리어 $G$ 에만 둔다.)
\item \textbf{집단별 꼬리 길이.} 장벽 $G$ 인 집단은 식~\eqref{eq:Lq},~\eqref{eq:keff} 로 자신의 꼬리
  길이 $L_{q}(G)\simeq\dfrac{|I|}{Q_\cell}\dfrac{h}{k_BT}\,e^{(G-\chi\mathcal A_j)/RT}$ 를 갖는다(\emph{활성화 지배}
  $k_j^{\eff}\simeq k_{j,\mathrm{act}}$ 형 — \emph{이하 \S\ref{sec:dist}$\cdot$\S\ref{sec:synth} 닫힌식 공통 scope};
  유한 공율속은 식~\eqref{eq:LqV} 의 감쇠인자와 \S\ref{sec:falsify} 경쟁원 배제로 다룬다). \emph{장벽이 길이로
  지수 변환}됨에 유의 — 작은 배리어 분산이 길이에서 지수적으로 확대된다.
\item \textbf{선형 중첩.} 관측되는 $\dd Q$ 는 각 집단의 기여 $\dd Q_G$ 의 \emph{선형 합}이다
  (독립 완화 가정에서 신호가 가산적). 각 집단 기여(식~\eqref{eq:tail}; 그 집단의 잔류 지연 $r_a(G)=r_j(q_a;G)$ 가
  진폭)를 $\rho_G$ 로 가중 적분하면(\S\ref{sec:stattools} 의 가중 평균, 식~\eqref{eq:wavg})
\end{enumerate}
\begin{equation}
\boxed{\;\begin{aligned}
\frac{\dd Q}{\dd V_n}\Big|_{\text{tail}}&=Q_j\!\int_0^\infty\!\rho_G(G)\,\frac{r_a(G)}{L_V(G)}\,
\exp\!\Big[-\frac{V_n-V_a}{L_V(G)}\Big]\,\dd G\quad\big(s(V_n-V_a)\ge0\big),\\[2pt]
&\qquad L_V(G)=\Big|\frac{\dd V_n}{\dd q}\Big|_{q_a}\!\!L_q(G),\qquad r_a(G)\equiv r_j(q_a;G).
\end{aligned}\;}
\label{eq:superpose}
\end{equation}
\textbf{중첩 결과의 해석.} 빠른 집단(작은 $G$, 작은 $L_q$)은 peak 근처에서 감쇠를 마치고, 느린 집단(큰 $G$)이 꼬리를 멀리까지
연장한다. ($\rho_G\to\delta(G-\bar G)$ 극한에서 적분이 sifting 으로 닫혀 식~\eqref{eq:simplefit} 의 단일 지수
($L_V=L_V(\bar G)$)로 \emph{정확히} 환원된다 — \S\ref{sec:synth} 와 정합.) \emph{용량 보존}: 일반 분포에서 꼬리 면적은
$Q_j\!\int\!\rho_G(G)\,r_a(G)\,\dd G=Q_j\overline{r_a}$ 이고, 포화 근사 $\overline{r_a}=1-\overline{\xi_j(q_a)}$ 로 미완 용량과
같아 전이 총 용량 $Q_j$ 를 넘지 않는다. 단일 지수의 \emph{연속 중첩}은 일반적으로
단일 지수가 아닌 \emph{stretched(늘어진)} 꼬리를 준다 \cite{lindsey1980,johnston2006,svare2000} —
\textbf{\S\ref{sec:lag} 가 예고한 ``실측 꼬리의 늘어짐''이 바로 이 분포 중첩의 귀결}이며, 단일 입자 단일지수
(\S\ref{sec:lag})는 분포가 좁은($\rho_G\to\delta$) 극한일 뿐이다.

\textbf{저온 꼬리 확장(분산 전파 — \S\ref{sec:stattools}).} 식~\eqref{eq:superpose} 의 $L_V(G)$ 는 $\ln(L_V/L_{V,0})=G/RT$(기준 길이 $L_{V,0}$ 가 구동항 $\chi\mathcal A_j$ 등 집단 공통 상수를 흡수해 무차원화;
$\because L_V\propto L_q\propto e^{(G-\chi\mathcal A_j)/RT}$ 이고 $\chi\mathcal A_j$ 는 상수 오프셋이라 비에서 탈락)이므로, 배리어 분산 $\sigma_G^2$ 가 $\mathrm{Var}[\ln(L_V/L_{V,0})]=\sigma_G^2/(RT)^2$(식~\eqref{eq:varprop} 을 $L_V\propto L_q$ 로 적용), 즉
표준편차 $\sigma_G/RT$ 로 전파된다. \textbf{저온일수록($RT\downarrow$) 길이 분산이 커져 꼬리가 더 늘어진다} — 고정된
배리어 분포가 온도를 통해 stretched 거동을 만드는 메커니즘이며, 이것이 관측의 ``저온 긴 꼬리''를 \emph{입자
분포까지 포함해} 설명한다.

\begin{rigorbox}
\textbf{길이 분포로의 좌표 변환(밀도 보존, 세 단계).} 식~\eqref{eq:superpose} 의 적분을 배리어 대신 길이 $L_q$
(전압축은 $L_V=|\dd V_n/\dd q|\,L_q$ 로 비례)로 옮기는 변환(\S\ref{sec:stattools} 의 변수 변환, 식~\eqref{eq:jacobian} 의 본 문제 형식 전개)을 명시한다. \emph{단계 1(밀도 보존)}: 집단의 분율은 좌표를 바꿔도 보존되므로 $A_L(L_q)\,\dd L_q=\rho_G(G)\,\dd G$.
\emph{단계 2(Jacobian)}: 역함수 $G(L_q)=\chi\mathcal A_j+RT\ln(L_q/L_0)$($L_0=|I|h/(Q_\cell k_BT)$)에서
$\dd G/\dd L_q=RT/L_q>0$ 이라 일대일이므로 $A_L(L_q)=\rho_G(G(L_q))\,RT/L_q$. \emph{단계 3(지지 범위)}: 배리어는
음수가 못 되므로($G\ge0$) 길이에 최단 하한 $L_{\min}=L_0\,e^{-\chi\mathcal A_j/RT}$ 가 생긴다($L_q<L_{\min}$ 에서
$A_L=0$). 본문에는 이 형식이 필요 없다 — 식~\eqref{eq:superpose} 의 ``$\rho_G$ 가중 합''만으로 충분하다.
\end{rigorbox}
\begin{boundbox}
\textbf{독립 완화 가정·비유일성.} 식~\eqref{eq:superpose} 는 입자 집단이 \emph{서로 독립}으로 완화한다고 본다
(평균장; 상경계 결합$\cdot$협동 효과 무시) \cite{funabiki1999}. 또 $\rho_G$ 가 배리어 $G$ \emph{하나}로 모든 집단
차이를 흡수한다는 것($r_a(G)\cdot L_V(G)$ 와의 상관이 평균장$\cdot$단일전위 $V_n$ 근사로 제한)이 본 절의 최강 가정이다.
또 stretched 꼬리에서 $\rho_G$ 로의 역산은
\emph{비유일}하므로($\rho_G$ 외 다른 메커니즘도 같은 꼬리를 낼 수 있다), $\rho_G$ 를 꼬리에서 역산하지 않고
독립 측정/모형한 $\rho_G$ 로 \emph{forward 예측}만 한다(\S\ref{sec:falsify}).
\end{boundbox}

% ====================================================================
\section{종합: 피팅 가능한 식과 3$\times$3 의존성}\label{sec:synth}
% ====================================================================
위 다섯 \emph{물리} 절(전하 보존 $\to$ 평형 $\to$ 지연 $\to$ 유효 배리어 $\to$ 분포; 통계 도구절
\S\ref{sec:stattools} 제외)을 묶으면 $\dd Q/\dd V$ 가 $(V;\,T,|I|,V_\drive)$ 의 닫힌 함수가 된다. peak 의
위치$\cdot$너비$\cdot$높이는 세 인자에 \emph{각각} 어떻게 반응하는가? 그리고 실무에 바로 쓸 가장 단순한 식은?

\subsection{닫힌 형태 — ``평형에서 지연을 뺀 것의 미분''}
\textbf{\S\ref{sec:dist} 의 꼬리를 상승부와 한 식으로 묶기.} \S\ref{sec:dist} 는 \emph{한 전이의 꼬리}(post-peak)만
배리어 분포 적분 식~\eqref{eq:superpose} 으로 줬다. 그러나 관측 peak 는 상승부(rise)와 꼬리가 \emph{한 곡선}이다.
이제 둘을 \emph{하나의 닫힌 식}으로 묶으면, \S\ref{sec:eqpeak} 의 평형 peak 가 rise 를, 식~\eqref{eq:superpose} 가
그 뒤 꼬리를 담당함이 한눈에 드러난다.

관측 peak 항은 식~\eqref{eq:dQdV} 의 $Q_j\,\dd\xi_j/\dd V_n$ 이다. 실제 전환율 $\xi_j$ 는 평형 목표 $\xi_{\eq,j}$ 를
지연만큼 못 따라가므로(\S\ref{sec:lag}), 그 모자란 양을 집단 평균한 \emph{가중 평균 지연}
\[
\bar r_j(q)\equiv\int_0^\infty\rho_G(G)\,r_j(q;G)\,\dd G,\qquad
r_j(q;G)=r_a(G)\,e^{-(q-q_a)/L_q(G)}\quad(\text{집단 }G\text{ 의 지연, 식~\eqref{eq:rsol}})
\]
을 빼서 $\xi_j=\xi_{\eq,j}-\bar r_j$ 로 쓴다. 즉 $r_j(q;G)$ 는 \S\ref{sec:lag} 가 \emph{한 집단}에 대해 구한 지연이고,
$\bar r_j$ 는 그것을 \S\ref{sec:dist} 의 분포 $\rho_G$ 로 평균(\S\ref{sec:stattools} 의 가중 평균)한 양이다.

\textbf{이 지연의 미분이 곧 식~\eqref{eq:superpose} 의 꼬리다(연결 확인).} $\xi_j=\xi_{\eq,j}-\bar r_j$ 를
식~\eqref{eq:dQdV} 에 넣으면 peak 항이 \emph{평형 미분} 빼기 \emph{지연 미분} 두 조각으로 갈린다. 꼬리를 만드는
지연 조각 $-Q_j\,\dd\bar r_j/\dd V_n$ 을 풀면, post-peak 에서 한 집단의 지연이 $\dd r_j/\dd q=
-(r_a(G)/L_q(G))\,e^{-(q-q_a)/L_q(G)}$ 로 감쇠하고, 연쇄율 $\dd/\dd V_n=(\dd q/\dd V_n)\,\dd/\dd q$ 와
$L_V(G)=|\dd V_n/\dd q|\,L_q(G)$ 를 넣으면
\[
-Q_j\frac{\dd\bar r_j}{\dd V_n}
=Q_j\!\int_0^\infty\!\rho_G(G)\,\frac{r_a(G)}{L_V(G)}\,e^{-(V_n-V_a)/L_V(G)}\,\dd G,
\]
즉 식~\eqref{eq:superpose} 의 꼬리가 \emph{정확히 재현}된다 — \S\ref{sec:dist} 의 분포 적분은 새 가정이 아니라
$\bar r_j$ 미분의 다른 표기였던 것이다. 그러므로 peak 전체는
\begin{equation}
\boxed{\;\frac{\dd Q}{\dd V_n}=C_\bg+Q_j\frac{\dd}{\dd V_n}\big(\xi_{\eq,j}-\bar r_j\big)
=\underbrace{C_\bg+Q_j\frac{\dd\xi_{\eq,j}}{\dd V_n}}_{\text{평형 peak (rise)}}
\;-\;\underbrace{Q_j\frac{\dd\bar r_j}{\dd V_n}}_{\text{지연이 깎는 꼬리}}\;.}
\label{eq:closed}
\end{equation}
이것이 종합식이다 — 모호한 결합 연산이 아니라 \emph{평형 전환율에서 지연을 뺀 것의 미분}이라는 \emph{정의된}
형태다. peak \emph{앞쪽}(상승부)에서는 지연이 작아 평형 항 $Q_j\dd\xi_{\eq,j}/\dd V_n$ 이 지배하고, peak
\emph{뒤쪽}에서는 $\dd\xi_{\eq,j}/\dd V_n\to0$ 이라 $-Q_j\dd\bar r_j/\dd V_n$ 의 지수 꼬리(식~\eqref{eq:superpose})가
지배한다. 빠른 kinetics 극한($\bar r_j\to0$)에서 평형 peak(\S\ref{sec:eqpeak})로 복귀하므로 식이 무결하다. 측정축은
$V_n\to V_\app=V_n+s_I|I|R_n$ 환산.

\subsection{3$\times$3 의존성 표 (실무 산출물)}
표의 \emph{행}은 peak 의 위치$\cdot$너비$\cdot$높이, \emph{열}은 세 인자($T\downarrow$$\cdot$$|I|\uparrow$$\cdot$$V_\drive\uparrow$)다.
각 칸은 그 인자가 그 특성을 바꾸는 \emph{경로} — \emph{어느 식의 어느 항}이 지배하고 부호가 어느 쪽인지 — 를 위
도출에서 직접 읽은 것이며($s=+1$ 방전 기준), 평형 항(\S\ref{sec:eqpeak})과 동역학 꼬리 항(\S\ref{sec:lag}$\cdot$\S\ref{sec:dist})이
\emph{반대로} 미는 칸(특히 너비$\cdot$높이의 $T$ 열)이 곧 관찰 해소의 핵심이다.
\begin{center}
\renewcommand{\arraystretch}{1.4}
\begin{tabular}{p{0.12\textwidth} p{0.26\textwidth} p{0.24\textwidth} p{0.25\textwidth}}
\toprule
 & \textbf{온도 $T\downarrow$} & \textbf{C-rate $|I|\uparrow$} & \textbf{구동 전위 $V_\drive\uparrow$} \\
\midrule
\textbf{위치} & $U_j(T)$ 이동, $\partial U_j/\partial T=\Delta S_j/(sF)$ \eqref{eq:dUdT} & 분극 $+s_I|I|R_n$ 이동 \eqref{eq:vapp} & 위치 직접 이동 없음 — $V_\peak=U_j$ 는 $V_\drive$ 무관 \eqref{eq:eqpeak}; \eqref{eq:LqV} 로 너비$\cdot$높이에만 \\
\textbf{너비} & 평형 $w_j=RT/F\downarrow$ \emph{좁힘} \eqref{eq:logistic}; 지연$\cdot$분포 꼬리 $\uparrow$ \emph{넓힘}($L_{V,j}\!\gtrsim\!w_j$ 면 우세) \eqref{eq:tailTC}\eqref{eq:superpose} & 꼬리 $L_{q,j}\propto|I|\uparrow$ 넓힘 \eqref{eq:tailTC} & 유효 배리어$\downarrow\Rightarrow L_{q,j}\downarrow$ 좁힘 \eqref{eq:LqV} \\
\textbf{높이} & 평형 $1/w_j\uparrow$, 지연$\cdot$분포가 낮춤(면적 보존) \eqref{eq:tailTC} & $|I|\uparrow\Rightarrow$ 낮춤 \eqref{eq:tailTC} & 배리어$\downarrow\Rightarrow$ 꼬리 짧아져 높임 \eqref{eq:LqV} \\
\bottomrule
\end{tabular}
\end{center}
\textbf{표 읽는 법(칸별 기전).} 표는 외울 것이 아니라 \emph{어느 식의 어느 항}에서 읽히는지를 보는 것이다.
\emph{위치 행} — 온도는 평형 중심 $U_j(T)$ 를 반응 엔트로피만큼 옮기고(식~\eqref{eq:dUdT}), C-rate 는 분극
$s_I|I|R_n$ 으로 \emph{apparent} 위치를 밀며(식~\eqref{eq:vapp}), 구동 전위는 평형 중심을 \emph{옮기지 못한다}
($V_\peak=U_j$ 는 $V_\drive$ 무관 — 전위는 평형 target 이 아니라 도달 \emph{속도}만 바꾸므로). \emph{너비 행} — 평형
폭 $w_j=RT/F$ 는 저온서 \emph{좁아}지지만(식~\eqref{eq:logistic}), 동역학 꼬리 $L_{V,j}$ 는 저온$\cdot$고율서
\emph{넓어}지고(식~\eqref{eq:tailTC}), 구동 전위는 유효 배리어를 낮춰 꼬리를 \emph{좁힌다}(식~\eqref{eq:LqV}).
\emph{높이 행} — 면적이 $\sim Q_j$ 로 보존되므로 너비와 \emph{반대}로 움직여, 평형은 저온서 높이고 꼬리는 낮춘다.
\textbf{이 표의 핵심은 너비$\cdot$높이의 $T$ 열에서 평형 항과 꼬리 항이 \emph{반대 부호}로 경쟁}한다는 점이다 —
평형만 보면 ``저온서 좁고 높다''인데, 꼬리가 이기면 관측은 ``저온서 넓고 낮다''로 뒤집힌다.

\textbf{핵심 관찰 해소.} 평형은 너비를 ``저온서 좁힘''으로 예측하나, 동역학 꼬리가 평형 폭보다 크면($L_{V,j}
\gtrsim\Delta V_{1/2}\sim w_j$; 저온$\cdot$유한 전류에서 일반적) 꼬리 넓힘이 우세해 \emph{저온서 넓고 비대칭}이
된다 — 관찰과 일치. (두 항의 \emph{크기} 우열은 데이터로 확인할 정량 문제이며, 본 장은 각 항의 \emph{부호}와
\emph{우세 조건} $L_{V,j}\gtrsim w_j$ 를 도출한다.) C-rate$\cdot$전위 방향도 가정이 아니라
위 식들에서 부호로 도출됐다.

\subsection{실무 피팅 근사식}
\textbf{우세 집단 극한.} $\rho_G$ 가 한 배리어 $\bar G$ 둘레에 좁게 몰리면($\rho_G\to\delta(G-\bar G)$),
식~\eqref{eq:superpose} 의 적분이 델타 sifting 으로 닫힌다($\int\delta(G-\bar G)f(G)\,\dd G=f(\bar G)$):
\begin{equation}
\boxed{\;L_{q,j}\simeq\frac{|I|}{Q_\cell}\frac{h}{k_BT}\,e^{(\bar G-\chi\mathcal A_j)/RT},\qquad
\frac{\dd Q}{\dd V_n}\Big|_{\text{tail}}\approx\frac{Q_j\,r_{a,j}}{L_{V,j}}\,e^{-(V_n-V_a)/L_{V,j}}\;,\;\;
L_{V,j}=\Big|\frac{\dd V_n}{\dd q}\Big|_{q_a}L_{q,j}\;.}
\label{eq:simplefit}
\end{equation}
즉 peak 는 \emph{평형 logistic 상승부 + 단일 지수 꼬리}이고, 꼬리는 특성 길이 $L_{V,j}$(전압 단위)의 지수 감쇠다
(꼬리는 식~\eqref{eq:tail} 처럼 post-peak \emph{단방향} $s(V_n-V_a)\ge0$ 에서만).
\textbf{단일지수 vs stretched 선택.} 분포가 좁으면($\rho_G\to\delta$) 이 단일지수가 정확하고(\S\ref{sec:lag} 환원),
넓어 꼬리가 \emph{휘면}(stretched) 식~\eqref{eq:superpose} 의 분포 적분으로 돌아가 $\rho_G$ 폭을 독립 입력으로
더한다 — \S\ref{sec:stattools}$\cdot$\S\ref{sec:dist} 의 분포는 \emph{버린 우회가 아니라} 단일지수가 깨지는
데이터에서 켜지는 항이다.
여기서 $r_{a,j}\equiv r_j(q_a)=\xi_{\eq,j}(q_a)-\xi_j(q_a)\in(0,1]$ 는 꼬리 시작점의 잔류 지연(전이가 꼬리로 넘기는
미완 분율)이며, 꼬리 면적은 $\int(\dd Q/\dd V_n)_{\text{tail}}\,\dd V_n=Q_j\,r_{a,j}$ 다. 정확히는 상승부가
$Q_j\,\xi_j(q_a)=Q_j[\xi_{\eq,j}(q_a)-r_{a,j}]$ 를, 꼬리가 $Q_j\,r_{a,j}$ 를 담아 합이 $Q_j\,\xi_{\eq,j}(q_a)$ 인데,
꼬리 영역의 \emph{포화 근사} $\xi_{\eq,j}(q_a)\simeq1$(식~\eqref{eq:Lq} 유도 전제) 하에서 상승부 $\simeq Q_j(1-r_{a,j})$,
총합 $\simeq Q_j$ 로 보존된다. \emph{이 인자 $r_{a,j}$ 없이 $Q_j/L_{V,j}$ 로 두면 꼬리 항만으로 $Q_j$ 가 소진되어 용량 보존이 깨진다.}
\begin{verifybox}
\textbf{Worked example — 한 전이의 peak 치수($Q_j=0.5\,Q_\cell$ 가정).} 평형 폭은 $25^\circ$C 에서
$w_j=RT/F=25.7$ mV 라 FWHM $=3.53\,w_j\approx91$ mV, 순 peak 높이 $=Q_j/(4w_j)=0.5\,Q_\cell/(4\times0.0257)
\approx4.9\,Q_\cell/\mathrm V$. 같은 전이가 $-15^\circ$C 에선 $w_j=22.2$ mV 로 좁아져 FWHM $\approx78$ mV, 높이
$\approx5.6\,Q_\cell/\mathrm V$ — \emph{평형만 보면 저온서 더 좁고 높다}. 그러나 동역학 꼬리 길이 $L_{q,j}\propto
e^{\Delta H_a/RT}$ 는 $\Delta H_a=40$ kJ/mol 이면 $25^\circ$C$\to-15^\circ$C 에서 약 $12$배 길어져(\S\ref{sec:lag}),
꼬리가 평형 폭을 넘어서면($L_{V,j}\gtrsim w_j$) 관측 peak 는 거꾸로 \emph{저온서 넓고$\cdot$낮고$\cdot$비대칭}이
된다 — 표의 너비$\cdot$높이 $T$ 열에서 평형과 꼬리가 반대로 미는 바로 그 칸이다. 이 한 예가 본 장 전체의
관찰(저온 긴 꼬리)을 수치로 압축한다.
\end{verifybox}

\textbf{식별 순서(비순환).} 모든 파라미터를 한 데이터에서 동시 자유 피팅하면 공선형이라 분리 불가다. 핵심 원리는
\emph{공선형을 한 겹씩 벗기는} 것이다 — 뒤 단계 파라미터와 꼬리에서 섞이는 양을 \emph{독립 실험}으로 \emph{먼저}
고정해 뒤 단계의 기지값으로 만들면, 각 단계에 미지수가 하나씩만 남아 순환이 끊긴다. 아래 \emph{순차} 제약이 그
벗김의 순서이며 — 각 단계는 \emph{무엇을 고정하고, 무엇이 분리되며, 공선형이 어떻게 해소되는가}로 읽는다.
\begin{enumerate}[label=\textbf{S\arabic*.},leftmargin=2.6em,itemsep=2pt]
\item \textbf{저율 OCV}($|I|\to0$, 꼬리가 사라진 평형). 고정$\to\{U_j,w_j,Q_j\}$. 분리 근거: 동역학 꼬리가
  없어 평형 peak(식~\eqref{eq:eqpeak})만 남으므로 위치$\cdot$너비$\cdot$면적이 곧 $U_j,w_j,Q_j$.
\item \textbf{펄스/전류 차단}. 고정$\to R_n$. 분리 근거: 전류 차단 직후 전압의 즉각 도약(IR-drop) $\Delta V=
  s_I|I|R_n$ 에서 $R_n$ 을 읽는다(GITT) \cite{weppner1977}. $R_n$ 과 $\chi$ 가 둘 다 꼬리를 늘려 공선형이므로
  $R_n$ 을 \emph{먼저} 묶는다.
\item \textbf{다전류/전위} ($\chi$ 를 \emph{Arrhenius 앞에서} 먼저). 고정$\to\chi$. 분리 근거: 식~\eqref{eq:LqV} 의
  $\partial\ln L_{q,j}/\partial V_\drive$ 기울기에서 $\chi$ 가 나온다(활성화 지배 극한 $-\chi sF/RT$; 공율속 유한이면
  식~\eqref{eq:LqV} 의 감쇠인자 $k_{j,\lim}/(k_{j,\mathrm{act}}+k_{j,\lim})$ 보정. $\mathcal A_j$ 는 S1$\cdot$S2 기지값). \textbf{$\chi$ 를 다음 단계(S4)보다 먼저 확정}해, S4 의 구동항 보정 $\chi\mathcal A_j$ 를 기지값으로
  만든다 — 이 순서가 ``$\chi$ 와 $\Delta H_a$ 가 둘 다 꼬리에 들어가 공선형''인 순환 의존을 끊어 비순환으로 식별한다.
  (실측: 고정 $T$, S1$\cdot$S2 로 환산한 동일 SOC$\cdot$tail window 에서 $\ln L_{q,j}$ 를 $V_\drive$ 로 회귀; 상세 절차는 피팅 장.)
\item \textbf{다온도 Arrhenius}. 고정$\to\Delta H_a$(순수, 기울기에서; $\Delta S_a$ 는 절편 $\Delta S_a/R+\text{const}$
  \emph{조합}으로만 — 절 끝에서 상술). \emph{보정의 대수 유도}: 식~\eqref{eq:simplefit}
  에 $\Delta G_a=\Delta H_a-T\Delta S_a$ 를 넣으면
  \begin{equation*}
  \begin{gathered}
  \ln L_{q,j}\simeq\ln\frac{T_*}{T}+\frac{\Delta H_a-T\Delta S_a-\chi\mathcal A_j}{RT},\\[2pt]
  T_*\equiv\frac{|I|h}{Q_\cell k_B}\ (\text{특성 온도, K — 로그 인자 무차원화}).
  \end{gathered}
  \end{equation*}
  $\ln(T_*/T)$ 의 $-\ln T$ 부분은 \emph{약한 온도 의존 보정항}, 엔트로피 항 $-\Delta S_a/R$(=$-T\Delta S_a/RT$)는
  \emph{상수 절편}($1/T$ 무관)이며, 구동항 $-\chi\mathcal A_j/RT$ 는 $\mathcal A_j(T)$ 때문에 $1/T$ 에 비선형으로
  섞인다. 좌변에서 $\ln(T_*/T)$ 와 구동항을 이항해
  \begin{equation}
  \boxed{\;\begin{aligned}
  &y(T)\equiv\ln\frac{T_*}{L_{q,j}\,T}-\frac{\chi\mathcal A_j(T)}{RT}\ \text{를}\ \tfrac1T\ \text{로 회귀}\\[2pt]
  &\quad\Rightarrow\ \text{기울기}=-\frac{\Delta H_a}{R}\ (\text{순수})
  \end{aligned}\;}
  \label{eq:arrhenius}
  \end{equation}
  (이항 후 남는 $1/T$ 계수는 $-\Delta H_a/R$ 뿐, 절편은 $\Delta S_a/R+\text{const}$). 분리 근거: prefactor$\cdot$
  엔트로피$\cdot$구동항을 보정해 순수 엔탈피만 기울기에 남긴다($\chi\mathcal A_j$ 는 \emph{앞 단계 S3} 의 기지 $\chi$ 와
  S1$\cdot$S2 로 계산되므로 순환이 없다).
\end{enumerate}
이로써 $\{U_j,w_j,Q_j,R_n,\Delta H_a,\chi\}$ 가 \emph{위 도출식만으로}(\emph{활성화 지배} 가정 하 — 유한 공율속이면
$k_{j,\lim}$ 이 추가 미지수라 \S\ref{sec:falsify} (i) 수송 배제로 닫는다) 데이터에서 분리 추출된다. \emph{단 $\Delta S_a$ 는}
Arrhenius 절편 $\Delta S_a/R+\text{const}$ 에서 \emph{현상론적 prefactor}(흡수된 $k_0\cdot$활성면적 등; \S\ref{sec:lag} 의
$k_0$ 주)와 \emph{결합}으로만 나오므로, prefactor 를 독립 고정하지 않으면 $\Delta S_a$ 단독값은 식별되지 않는다(절편
조합으로만 결정). $\rho_G$(분포 폭)는 꼬리가 단일 지수에서 벗어나(stretched) 휠 때 독립 입력으로 더한다(역산 금지).
\begin{boundbox}
\textbf{본 장 범위와 식별 교란 요인.} (i) 본 장은 \emph{단일 방향}(방전)$\cdot$단일 전이를 다룬다 —
lithiation/delithiation 간 OCV hysteresis 는 peak 위치의 방향 비대칭을 따로 만들므로 본 장 밖이다(충전 branch
확장에서). (ii) 여러 전이가 겹치는 영역(관측의 ``다음 peak 와 겹침'')은 \S\ref{sec:overlap} 에서 전이 $j$ 별
합산$\cdot$OCV-anchored 분해로 본격 다룬다. (iii) 배경 용량 $C_\bg$(식~\eqref{eq:dQdV})은 꼬리와 공선형이
될 수 있어, S1 에서 저율 OCV 로 함께 고정한 뒤 그 위에서 꼬리를 회귀한다. (iv) 본 장 전위는 \emph{반쪽전지}(상대전극
기준 내부 $V_n$) 척도이며 ICA 도 $\dd Q/\dd V_n$ 이다 — full-cell $V_\cell=V_p-V_n$ 의 $\dd Q/\dd V_\cell$ 는 양극
기여가 섞이므로, 본 장 결과를 full-cell 데이터에 적용하려면 양극 분리(반쪽전지 기준 환원)가 선행돼야 한다.
\end{boundbox}

% ====================================================================
\section{다중 전이 겹침과 합산 분해}\label{sec:overlap}
% ====================================================================
\S\ref{sec:synth} 까지 \emph{단일 전이}의 peak 를 닫았다. 그러나 관측의 한 축은 ``C-rate 가 커지면 peak 가
\emph{겹친다}''였고, 흑연은 전이가 여럿이다($N_p\approx3\sim4$). 그 단독 곡선들이 겹칠 때 관측 ICA 는 어떻게
합쳐지며, 저율 정보로 그 합을 어떻게 \emph{분해$\cdot$예측}하는가?

\textbf{총 ICA 의 선형 중첩.} 전하 보존식~\eqref{eq:charge} 은 처음부터 $\sum_j Q_j\xi_j$ 형이었고, 그 미분
~\eqref{eq:dQdV} 이 곧 총 ICA 다:
\begin{equation}
\boxed{\;\frac{\dd Q}{\dd V_n}=C_\bg(V_n,T)+\sum_{j=1}^{N_p}Q_j\,\frac{\dd\xi_j}{\dd V_n}\;}\qquad(\text{식~\eqref{eq:dQdV} 의 다전이 형})
\label{eq:total}
\end{equation}
각 항 $Q_j\,\dd\xi_j/\dd V_n$ 은 식~\eqref{eq:closed} 의 (평형 rise $-$ 지연 꼬리) 구조다. 전이 집단이 서로
독립이면(\S\ref{sec:dist} 평균장) 신호가 가산적이라 관측 ICA 는 \emph{각 전이 단독 곡선의 합}이다 — 새 가정이
아니라 전하 보존의 직접 미분이다.

\textbf{겹침의 정량 기준.} 인접 전이 $j,j{+}1$ 의 단독 곡선이 겹치는 조건은 둘이다: (i) \emph{평형} peak
폭으로도 가까우면 — 중심 차가 두 평형 폭 평균 이하 $|U_{j+1}-U_j|\lesssim(\Delta V_{1/2,j}+\Delta V_{1/2,j+1})/2$
(각 $\simeq3.53\,w$; 식~\eqref{eq:eqpeak}); (ii) \emph{동역학} 꼬리가 옆 peak 까지 닿으면 — 꼬리 길이가 중심 차 이상
$L_{V,j}\gtrsim|U_{j+1}-U_j|$(식~\eqref{eq:tail}; 꼬리는 단방향이라 $j$ 의 꼬리는 방전 방향 \emph{다음} 전이 쪽으로
뻗는다). \textbf{(ii)가 C-rate$\cdot$저온의 겹침이다}: 식~
\eqref{eq:tailTC} 로 $|I|\uparrow$ 또는 $T\downarrow$ 면 $L_{V,j}\uparrow$ 라, 저율서 분리됐던 peak 가 고율서 조건
(ii)를 넘겨 융합한다. 융합 시작 전류$\cdot$온도는 위 부등식의 등호로 추정된다.
\begin{verifybox}
\textbf{Worked example — 두 peak 의 융합 개시.} 인접 stage 전이가 $\Delta U=|U_{j+1}-U_j|\approx100$ mV 떨어져
있고 평형 폭 $w_j=25.7$ mV(FWHM $\approx91$ mV)면, 저율 평형서는 평균 FWHM($\approx91$ mV)이 간격($100$ mV)보다
작아 \emph{겨우 분리}된다. 전류를 키워 꼬리 $L_{V,j}$ 가 $\sim100$ mV 로 자라면(조건 (ii)) $j$ 의 꼬리가 $j{+}1$
peak 에 닿아 \emph{융합이 시작}된다. $L_{V,j}\propto|I|\,e^{\Delta H_a/RT}$ 라 저온일수록 더 \emph{낮은} 전류에서
융합이 개시된다 — 관측의 ``저온$\cdot$고율서 peak 가 겹친다''가 이 부등식의 등호로 정량화된다. 일단 융합하면 두
전이의 면적이 한 봉우리에 섞여, 저율 anchor 없이는 $Q_j$ 를 되돌릴 수 없다(아래).
\end{verifybox}

\textbf{OCV-anchored 합산 분해$\cdot$예측(절차).} 융합한 고율 데이터에서 $j$별 기여를 \emph{역산}하면 비유일
이다(\S\ref{sec:dist}). 대신 저율 정보를 anchor 로 \emph{forward} 예측한다:
\begin{enumerate}[label=\textbf{S$_\arabic*'$.},leftmargin=2.8em,itemsep=2pt]
\item \textbf{저율 OCV 로 분리 peak 에서 $\{U_j,w_j,Q_j\}_{j=1}^{N_p}$ 고정}(식~\eqref{eq:eqpeak}; \S\ref{sec:synth}
  의 S1 과 동일). 저율서는 꼬리가 짧아 peak 가 분리되므로 전이별 평형 파라미터가 식별된다.
\item \textbf{각 전이의 고율 단독 곡선 생성}: 고정 $\{U_j,w_j,Q_j\}$ 에 동역학 $\{\Delta H_{a,j},\chi\}$ 와 Arrhenius
  \emph{절편 조합값}(\S\ref{sec:synth} S3$\cdot$S4 — $\Delta S_a$ 단독이 아니라 흡수 prefactor 와 결합된 절편이
  $L_{q,j}$ 절대값을 정함)을 넣어 $j$별 평형 rise $+$ 지수 꼬리(식~\eqref{eq:simplefit})를 그린다.
\item \textbf{합산} 식~\eqref{eq:total} 로 \emph{융합 형상을 예측}한다($\sum_j$).
\item \textbf{데이터 대조}: 예측 융합 곡선을 고율 측정과 맞춰 동역학 파라미터를 정련한다 — $\{U_j,w_j,Q_j\}$ 는
  S$_1'$ 고정값 유지(자유화하면 공선형).
\end{enumerate}
\textbf{분리 가능성의 한계.} peak 가 완전히 융합하면 $j$별 \emph{면적}($Q_j$) 분리는 고율 데이터만으로는 불가능
하다(\S\ref{sec:synth} 공선형 논리의 겹침판). \emph{저율 OCV anchor 가 있을 때만} 융합 peak 의 합산 정보를
전이별로 되돌릴 수 있다 — 그래서 본 절은 ``저율로 고정, 고율로 예측''의 forward 방향만 허용한다(역산 금지,
\S\ref{sec:dist} 비유일성과 정합).

% ====================================================================
\section{충방전 분기와 히스테리시스 gap $\Delta U_j^\hys$}\label{sec:hys_branch}
% ====================================================================
지금까지(평형 peak $+$ 정규용액 상분리 $+$ 동역학)는 \emph{방전 한 방향}의 $\dd Q/\dd V$ 를 닫았다. 이제 같은
전극을 충전과 방전으로 오갈 때 — 같은 SOC 에서 전위가 갈리는 \emph{충방전 히스테리시스} — 로 넓힌다. 그 뿌리는
\S\ref{sec:regsol} 의 상분리다: 충$\cdot$방전이 서로 다른 준안정 분기(spinodal 과주행)를 따라가기 때문.

\subsection{비단조 평형 전위 $V_{\eq,j}(\xi)$}
\S\ref{sec:regsol} 의 상분리를 \emph{전위}의 언어로 옮긴다. 평형 조건 식~\eqref{eq:eqcond}
$\mu_\mathrm{Li}(\theta_{\eq,j})=\mu^0-sF(V_n-U_j)$ 에 화학퍼텐셜 식~\eqref{eq:mu}(이제 $\Omega_j\ne0$)를 넣고
$\theta=1-\xi$ 로 바꾸면 — $\ln[\theta/(1-\theta)]=-\ln[\xi/(1-\xi)]$, $(1-2\theta)=-(1-2\xi)$ 라 두 항 부호가 함께
뒤집혀 —
\begin{equation}
V_{\eq,j}(\xi)=U_j+\frac{RT}{sF}\ln\frac{\xi}{1-\xi}+\frac{\Omega_j}{sF}(1-2\xi).
\label{eq:hys_Veq}
\end{equation}
$\Omega_j=0$ 이면 둘째 항만 남아 앞 절 logistic 식~\eqref{eq:logistic} 의 역함수(단조)다. $\Omega_j$ 항이
\emph{되돌림}을 더한다. 기울기는
\begin{equation}
\frac{\dd V_{\eq,j}}{\dd\xi}=\frac{1}{sF}\Big[\frac{RT}{\xi(1-\xi)}-2\Omega_j\Big]=\frac{1}{sF}\,g_j''(\xi),
\label{eq:hys_slope}
\end{equation}
즉 \emph{전위 곡선의 기울기가 곧 \S\ref{sec:regsol} 의 자유에너지 곡률}(식~\eqref{eq:regsol_gpp})이다 —
$V_{\eq,j}$ 기울기 $0$ 인 곳이 spinodal($g_j''=0$). $\Omega_j>2RT$ 면 $\xi=\tfrac12$($RT/[\xi(1-\xi)]=4RT$) 둘레서
음의 기울기(역학적 불안정)다. 기울기 $0$ 인 spinodal 진행률은 식~\eqref{eq:hys_slope}$=0$ 에서
$\xi(1-\xi)=RT/(2\Omega_j)$, 곧
\begin{equation}
\boxed{\;\xi_{s,j}^\pm=\tfrac12\pm\tfrac12\sqrt{1-\tfrac{2RT}{\Omega_j}}\;,\qquad(\Omega_j>2RT).\;}
\label{eq:hys_spinodal}
\end{equation}
\begin{verifybox}
\textbf{임계$\cdot$환원.} 실수 spinodal 은 $\Omega_j>2RT$($T<T_{c,j}=\Omega_j/2R$, \S\ref{sec:regsol})에서만. $\Omega_j
\to2RT^+$ 면 $\xi_{s,j}^\pm\to\tfrac12$, $\Omega_j\to\infty$ 면 $\to0,1$. $\Omega_j\le2RT$ 면 $V_{\eq,j}$ 단조라
spinodal$\cdot$히스 없음 — 앞 절 단일 peak.
\end{verifybox}

\subsection{충방전 분기와 gap}
비단조 $V_{\eq,j}$ 는 $\xi_{s,j}^-$ 서 \emph{극대}, 불안정 구간서 하강, $\xi_{s,j}^+$ 서 \emph{극소}, 다시 상승한다.
한 전이는 불안정 구간을 균질하게 못 지나 \S\ref{sec:regsol} 의 준안정 분기를 따른다. \textbf{방전}(탈리튬화, $\xi$
증가)은 상승 분기를 극대 $\xi_{s,j}^-$ 까지, \textbf{충전}(리튬화, $\xi$ 감소)은 반대 끝에서 극소 $\xi_{s,j}^+$ 까지
과주행한 뒤 상분리한다. 분기의 유효 중심 전위는 그 극값이고
\begin{equation}
U_j^\dis=V_{\eq,j}(\xi_{s,j}^-),\qquad U_j^\chg=V_{\eq,j}(\xi_{s,j}^+),
\label{eq:hys_branchU}
\end{equation}
극대$>$극소라 $U_j^\dis>U_j^\chg$ — 방전 전위가 충전보다 높다는 관찰과 부호가 맞는다. gap 은 두 극값 차다.
$\xi_{s,j}^\pm=\tfrac12\pm\tfrac12 u_j$($u_j\equiv\sqrt{1-2RT/\Omega_j}$)를 식~\eqref{eq:hys_Veq} 에 넣어 빼면 —
방전 끝점서 $\tfrac{\xi}{1-\xi}=\tfrac{1-u_j}{1+u_j}$, 충전서 $\tfrac{1+u_j}{1-u_j}$, $(1-2\xi)=\pm u_j$ — 로그 항
차 $2\ln\tfrac{1-u_j}{1+u_j}=-4\,\mathrm{artanh}\,u_j$, $\Omega_j$ 항 차 $2u_j$ 라
\begin{equation}
\boxed{\;\Delta U_j^\hys=U_j^\dis-U_j^\chg=\frac{2}{sF}\big[\Omega_j u_j-2RT\,\mathrm{artanh}\,u_j\big],\quad
u_j=\sqrt{1-\tfrac{2RT}{\Omega_j}}\;.}
\label{eq:hys_dU}
\end{equation}
이로써 히스테리시스 gap 이 \emph{정규용액 $\Omega_j$ 와 온도에서 유도}됐다 — 임의 피팅 상수가 아니다.
\begin{verifybox}
\textbf{극한$\cdot$수치.} $\Omega_j\to2RT^+$ 면 $\Omega_j u_j-2RT\,\mathrm{artanh}\,u_j\to\tfrac43RT u_j^3$ 라
$\Delta U_j^\hys\to(8RT/3sF)u_j^3\to0^+$(임계서 매끄럽게 소멸). $\Omega_j=4RT$ 면 $u_j=\sqrt{0.5}\approx0.707$,
$\mathrm{artanh}(0.707)\approx0.881$, $\Delta U_j^\hys\approx(2/sF)(4{\cdot}0.707-2{\cdot}0.881)RT\approx2.1\,RT/(sF)
\approx54$ mV($25^\circ$C). 차원: $\Omega,RT$ J/mol, $sF$ C/mol $\Rightarrow$ V.
\end{verifybox}
\begin{boundbox}
\textbf{spinodal 상한과 축소 $\gamma_j$.} 식~\eqref{eq:hys_branchU} 는 각 분기가 spinodal 까지 과주행하는 \emph{상한}
이다. 실제로는 핵생성이 spinodal 전에 시작되고 다입자 평활화(\S\ref{sec:regsol} Dreyer)가 갈림을 좁히므로, 실측
gap 은 그 아래다. 한 자유도 $\gamma_j\in[0,1]$ 로 흡수해 \emph{유효 분기 중심}을
\begin{equation}
U_j^{\dis/\chg}=U_j\pm\tfrac12\gamma_j\Delta U_j^\hys
\label{eq:hys_center}
\end{equation}
로 둔다($\gamma_j=1$ spinodal 상한, $\gamma_j\to0$ binodal=가역). $\Omega_j$ 와 $\gamma_j$ 한 양씩이면 충분하다.
\end{boundbox}

% ====================================================================
\section{종합 모델식과 피팅$\cdot$예측 — 한 줄로}\label{sec:master}
% ====================================================================
앞 절들이 \emph{각 절은 그 식의 한 항}이라 했다. 이제 그 항들을 하나로 모은다. 관측 $\dd Q/\dd V$ 는
\emph{(배경) $+$ 전이별 (평형 상승부 $+$ 동역학 꼬리)}의 합이며, 이것이 본 장이 도달하는 \emph{단일 피팅식}이다.

\subsection{최종 모델식}
측정 전위 $V_\app$ 를 내부 전위 $V_n=V_\app-s_I|I|R_n$(식~\eqref{eq:vapp})으로 환산한 뒤, 전이 $j=1,\dots,N_p$ 에 대해
\begin{equation}
\boxed{\;\frac{\dd Q}{\dd V_\app}=C_\bg(V_n,T)+\sum_{j=1}^{N_p}Q_j\bigg[\underbrace{\frac{\xi_{\eq,j}(1-\xi_{\eq,j})}{w_j}}_{\text{평형 상승부}}
\;+\;\underbrace{\frac{r_{a,j}}{L_{V,j}}\,e^{-(V_n-V_{a,j})/L_{V,j}}}_{\substack{\text{동역학 꼬리}\\ s(V_n-V_{a,j})\ge0}}\bigg]\;.}
\label{eq:master}
\end{equation}
각 항은 앞 절의 \emph{닫힌} 결과를 그대로 부른다:
\begin{align*}
\xi_{\eq,j}&=\big[1+e^{-s(V_n-U_j)/w_j}\big]^{-1} &&\text{(식~\eqref{eq:logistic}; 평형 상승부 → \S\ref{sec:eqpeak})},\\
L_{V,j}&=\Big|\tfrac{\dd V_n}{\dd q}\Big|_{q_a}L_{q,j},\quad
L_{q,j}=\tfrac{|I|}{Q_\cell}\tfrac{h}{k_BT}\,e^{(\Delta H_{a,j}-T\Delta S_{a,j}-\chi\mathcal A_j)/RT}
&&\text{(식~\eqref{eq:Lq}\eqref{eq:keff}; 꼬리 길이 → \S\ref{sec:lag}\S\ref{sec:barrier})},\\
\mathcal A_j&=sF(V_n-U_j),\quad r_{a,j}=\xi_{\eq,j}(q_a)-\xi_j(q_a)&&\text{(식~\eqref{eq:affinity}; 구동력$\cdot$잔류 지연)}.
\end{align*}
$V_{a,j}=V_n(q_{a,j})$=전이 $j$ 의 꼬리 시작 전위. 분포가 좁으면($\rho_G\to\delta$) 이 단일지수 꼬리가 정확하고,
넓어 꼬리가 휘면(stretched) 그 꼬리 항을 식~\eqref{eq:superpose} 의 $\rho_G$ 적분으로 바꾼다(\S\ref{sec:dist}).
\emph{이 한 식이 peak 의 위치$\cdot$너비$\cdot$높이와 그 온도$\cdot$C-rate$\cdot$전위 의존(3$\times$3, \S\ref{sec:synth})을 모두 담는다.}

\subsection{피팅 파라미터 — 전이당 무엇을, 어느 데이터로}
\begin{center}
\renewcommand{\arraystretch}{1.3}
\begin{tabular}{p{0.10\textwidth} p{0.30\textwidth} p{0.34\textwidth} p{0.14\textwidth}}
\toprule
파라미터 & 의미 & 닫는 데이터 & 식 \\
\midrule
$U_j$ & 평형 중심 전위 & 저율 OCV peak \emph{위치} & \eqref{eq:eqpeak} \\
$w_j$ & 평형 폭(이상 $RT/F$) & 저율 OCV peak \emph{너비} & \eqref{eq:logistic} \\
$Q_j$ & 전이 용량 & 저율 OCV peak \emph{면적} & \eqref{eq:eqpeak} \\
$R_n$ & 분극 저항 & 전류 차단 IR-drop(GITT) & \eqref{eq:vapp} \\
$\chi$ & 전달 계수 & 다율 $\partial\ln L_{q,j}/\partial V$ & \eqref{eq:LqV} \\
$\Delta H_{a,j}$ & 활성화 엔탈피 & 다온도 Arrhenius 기울기 & \eqref{eq:arrhenius} \\
$C_\bg$ & 배경 미분용량 & 저율 OCV 기저선 & \eqref{eq:dQdV} \\
($\rho_G$ 폭) & 배리어 분포 폭(선택) & 꼬리 stretched 휨 & \eqref{eq:superpose} \\
\bottomrule
\end{tabular}
\end{center}
즉 \emph{전이당 핵심 $\{U_j,w_j,Q_j,\Delta H_{a,j}\}$ 와 전이 공통 $\{\chi,R_n,C_\bg\}$}; 꼬리가 휠 때만 $\rho_G$ 폭을 더한다.

\subsection{데이터$\to$파라미터(비순환)$\to$예측}
식별은 \S\ref{sec:synth} 의 \textbf{S1$\to$S4} 순서다 — 각 단계가 앞 단계 기지값을 써 미지수를 하나씩만 남겨 공선형을 끊는다:
\textbf{S1} 저율 OCV($|I|\!\to\!0$)$\to\{U_j,w_j,Q_j,C_\bg\}$; \textbf{S2} 전류 차단$\to R_n$; \textbf{S3} 다율
$\to\chi$; \textbf{S4} 다온도 Arrhenius$\to\Delta H_{a,j}$. 여러 전이가 겹치면 식~\eqref{eq:total} 의 합산으로
저율-anchored \emph{forward} 예측한다(역산 금지 — \S\ref{sec:overlap}).
\begin{keybox}
\textbf{이 식 하나로 끝난다.} 저율 OCV $+$ 전류 차단 $+$ 다율$\cdot$다온도(예 OCV$\cdot$0.1C$\cdot$0.2C $\times$
15/23/45$^\circ$C)로 위 파라미터를 닫으면, \emph{같은 식~\eqref{eq:master} 에 임의의} $(T,|I|,V_\drive)$ 를 넣어
그 조건의 방전 $\dd Q/\dd V$ 를 \emph{예측}한다. 피팅은 식 하나, 파라미터는 전이당 소수, 데이터원은 분리돼 단계적.
\end{keybox}

\subsection{데이터에서 예측까지 — 한 walkthrough}
추상적 절차를 한 흐름으로 굳힌다. 한 셀의 흑연 음극을 예로:
\begin{enumerate}[label=\textbf{\arabic*.},leftmargin=2.0em,itemsep=1pt]
\item \textbf{준평형}($0.05$C 또는 GITT) 방전 $\dd Q/\dd V$ 에서 분리된 $N_p$ 개 peak 의 위치$\cdot$너비$\cdot$면적을
  읽어 $\{U_j,w_j,Q_j\}$ 와 기저선 $C_\bg$ 를 고정한다(S1).
\item 같은 SOC 들에서 전류 차단 IR-drop 으로 $R_n$ 을 얻는다(S2).
\item $0.1$C$\cdot$$0.2$C 두 율의 동일 tail window 에서 $\ln L_{q,j}$ 를 $V_\drive$ 로 회귀, 기울기에서 $\chi$(S3).
\item $15/23/45^\circ$C 의 $\ln L_{q,j}$ 를 구동항 보정해 $1/T$ 로 회귀, 기울기에서 $\Delta H_{a,j}$(S4).
\end{enumerate}
이제 식~\eqref{eq:master} 에 이 파라미터와 \emph{원하는} $(T,|I|)$ 를 넣으면 그 조건의 곡선이 나온다. \textbf{검증
시험}: 피팅에 \emph{쓰지 않은} 조건(예 $0.5$C$\cdot$$5^\circ$C)을 예측하고 그 측정과 대조한다 — 맞으면 모델이
서고, 빗나가면 \S\ref{sec:falsify} 의 경쟁 원인(확산$\cdot$전하전달)이 살아 있다는 신호다. \emph{피팅이 아니라
외삽 예측이 진짜 시험}이다.

% ====================================================================
\section{검증과 반증(falsification)}\label{sec:falsify}
% ====================================================================
지금까지 peak 의 모양과 그 인자 의존을 도출하고 겹침까지 예측했다. 좋은 이론은 \emph{반증 가능}해야 하므로,
마지막으로 묻는다: ``꼬리는 전위가 낮춘 유효 배리어의 동역학 산물''이라는 본 장의 주장은 \emph{어떻게 틀릴 수
있는가}? 그리고 같은 꼬리를 내는 \emph{경쟁 원인}과 어떻게 구별하는가?

\textbf{반증 가능한 예측.} 식~\eqref{eq:arrhenius} 에서, 꼬리 길이의 Arrhenius 기울기를 구동항으로 보정한 $y(T)$
의 기울기는 \emph{순수} $-\Delta H_a/R$ 이어야 하며 \textbf{구동 전위에 무관}해야 한다. 보정 전 유효 기울기는(해당
온도창서 $\mathcal A_j$ 거의 상수 $+$ \emph{활성화 지배} $k_{j,\mathrm{act}}\ll k_{j,\lim}$ 의 \emph{국소} 근사로 — 후자는
아래 (i) 수송 진단으로 \emph{검증}, 미수행 시 가정) $-(\Delta H_a-\chi\mathcal A_j)/R$ 라 \emph{전위에 의존}한다(일반적으로 $\mathcal A_j(T)$
비선형성$\cdot$유한 공율속의 직렬 율속 항이 기울기를 추가 왜곡하나, 핵심은 \emph{전위 의존 자체}다) — \textbf{이 전위
의존이 ``전위가 배리어를 낮춘다''의 반증 가능한 지문}이다. 보정 후에도 전위 의존이 남으면 본 설명은 기각된다.

\textbf{경쟁 꼬리원의 배제.} 꼬리 길이 $L_{q,j}\propto|I|$ 와 ``전위$\uparrow\Rightarrow$꼬리
$\downarrow$''는 \emph{단순 확산$\cdot$접촉저항$\cdot$전하전달 과전위}도 공유한다. 따라서 그것만으로
barrier-lowering 에 귀속하면 오귀속이다. 두 분리가 \emph{모두} 필요하다:
\begin{enumerate}[label=(\roman*),leftmargin=2.0em,itemsep=2pt]
\item \textbf{확산 배제.} 고체상 확산이 율속이고 확산계수가 점유율 의존이면 transport-limited 유효 엔탈피도 전위
  의존이다. GITT 완화 transient 형상($\sqrt t$ Cottrell=반무한 확산 vs 지수=계면 율속)으로 이를 \emph{독립 판정}해
  배제한 뒤에만 위 지문이 유효하다(이 형상 판정은 표준 진단으로 본 장 이론의 \emph{입력}이다) \cite{weppner1977}.
\item \textbf{전하전달 분리.} charge-transfer 과전위 $\eta_\mathrm{ct}$ 도 전위 지수의존이라 단일 전극서 $\chi$ 와
  공선형이다. GITT/전류 차단(또는 EIS)으로 $R_\mathrm{ct}$ 를 독립 분리하는 표준 전기화학 진단 뒤에만 $\chi$ 의
  ``배리어 낮춤'' 귀속이 성립한다 \cite{weppner1977,dubarry2012}.
\end{enumerate}
세 후보는 \emph{꼬리의 거시 거동}($\propto|I|$, 전위$\uparrow\Rightarrow$꼬리$\downarrow$)을 \emph{공유}하므로
거시 관측만으론 못 가른다 — 갈라 주는 것은 서로 다른 \emph{독립 진단}이다:
\begin{center}
\renewcommand{\arraystretch}{1.35}
\begin{tabular}{p{0.26\textwidth} p{0.30\textwidth} p{0.32\textwidth}}
\toprule
원인 & 전위 의존의 출처 & 갈라 주는 독립 진단 \\
\midrule
배리어 낮춤(본 장) & 유효 배리어 $-\chi\mathcal A_j$ & 구동항 보정 후 Arrhenius 기울기가 \emph{전위 무관} 하면 성립(식~\eqref{eq:arrhenius}) \\
고체상 확산 & $D(\theta)$ 의 점유율 의존 & GITT 완화가 $\sqrt t$(Cottrell) 형 — 계면 율속의 지수형과 다름 \\
전하전달 과전위 & $\eta_\mathrm{ct}\propto\ln i$ & EIS 의 $R_\mathrm{ct}$ 반원(또는 전류 차단)으로 독립 분리 \\
\bottomrule
\end{tabular}
\end{center}
즉 \textbf{반증의 칼날은 ``보정 후에도 남는 전위 의존''}이다 — 본 장의 barrier-lowering 이 옳다면 그것은 \emph{0}
이어야 하고, $\sqrt t$ 완화나 $R_\mathrm{ct}$ 반원이 보이면 그만큼을 먼저 떼어내야 한다. 이 분리 없이 단일
데이터셋에서 꼬리를 barrier-spectrum 으로 귀속하는 것은 금지한다.

% ====================================================================
\begin{thebibliography}{99}
\bibitem{dahn1991} J. R. Dahn, ``Phase diagram of Li$_x$C$_6$,'' \emph{Phys. Rev. B} \textbf{44}, 9170 (1991). DOI: 10.1103/PhysRevB.44.9170.
\bibitem{ohzuku1993} T. Ohzuku, Y. Iwakoshi, K. Sawai, ``Formation of Lithium-Graphite Intercalation Compounds in Nonaqueous Electrolytes,'' \emph{J. Electrochem. Soc.} \textbf{140}, 2490 (1993). DOI: 10.1149/1.2220849.
\bibitem{funabiki1999} A. Funabiki, M. Inaba, Z. Ogumi et al., ``Stage Transformation of Lithium-Graphite Intercalation Compounds Caused by Electrochemical Lithium Intercalation,'' \emph{J. Electrochem. Soc.} \textbf{146}, 2443 (1999). DOI: 10.1149/1.1391953.
\bibitem{dreyer2010} W. Dreyer \emph{et al.}, ``The thermodynamic origin of hysteresis in insertion batteries,'' \emph{Nature Materials} \textbf{9}, 448 (2010). DOI: 10.1038/nmat2730.
\bibitem{sethna1993} J. P. Sethna \emph{et al.}, ``Hysteresis and hierarchies: Dynamics of disorder-driven first-order phase transformations,'' \emph{Phys. Rev. Lett.} \textbf{70}, 3347 (1993). DOI: 10.1103/PhysRevLett.70.3347.
\bibitem{mckinnon1983} W. R. McKinnon, R. R. Haering, ``Physical Mechanisms of Intercalation,'' in \emph{Modern Aspects of Electrochemistry} \textbf{15}, 235 (1983). DOI: 10.1007/978-1-4615-7461-3\_4.
\bibitem{bazant2013} M. Z. Bazant, ``Theory of Chemical Kinetics and Charge Transfer based on Nonequilibrium Thermodynamics,'' \emph{Acc. Chem. Res.} \textbf{46}, 1144 (2013). DOI: 10.1021/ar300145c.
\bibitem{eyring1935} H. Eyring, ``The Activated Complex in Chemical Reactions,'' \emph{J. Chem. Phys.} \textbf{3}, 107 (1935). DOI: 10.1063/1.1749604.
\bibitem{marcus1956} R. A. Marcus, ``On the Theory of Oxidation-Reduction Reactions Involving Electron Transfer. I,'' \emph{J. Chem. Phys.} \textbf{24}, 966 (1956). DOI: 10.1063/1.1742723.
\bibitem{evanspolanyi1938} M. G. Evans, M. Polanyi, ``Inertia and driving force of chemical reactions,'' \emph{Trans. Faraday Soc.} \textbf{34}, 11 (1938). DOI: 10.1039/TF9383400011.
\bibitem{bardfaulkner} A. J. Bard, L. R. Faulkner, \emph{Electrochemical Methods: Fundamentals and Applications}, 2nd ed., Wiley (2001). ISBN 978-0-471-04372-0.
\bibitem{reynier2004} Y. Reynier, R. Yazami, B. Fultz, ``Thermodynamics of Lithium Intercalation into Graphites and Disordered Carbons,'' \emph{J. Electrochem. Soc.} \textbf{151}, A422 (2004). DOI: 10.1149/1.1646152.
\bibitem{lindsey1980} C. P. Lindsey, G. D. Patterson, ``Detailed comparison of the Williams--Watts and Cole--Davidson functions,'' \emph{J. Chem. Phys.} \textbf{73}, 3348 (1980). DOI: 10.1063/1.440530.
\bibitem{johnston2006} D. C. Johnston, ``Stretched exponential relaxation arising from a continuous sum of exponential decays,'' \emph{Phys. Rev. B} \textbf{74}, 184430 (2006). DOI: 10.1103/PhysRevB.74.184430.
\bibitem{svare2000} I. Svare, S. W. Martin, F. Borsa, ``Stretched exponentials with $T$-dependent exponents from fixed distributions of energy barriers for relaxation times in fast-ion conductors,'' \emph{Phys. Rev. B} \textbf{61}, 228 (2000). DOI: 10.1103/PhysRevB.61.228.
\bibitem{weppner1977} W. Weppner, R. A. Huggins, ``Determination of the Kinetic Parameters of Mixed-Conducting Electrodes (GITT),'' \emph{J. Electrochem. Soc.} \textbf{124}, 1569 (1977). DOI: 10.1149/1.2133112.
\bibitem{dubarry2012} M. Dubarry, C. Truchot, B. Y. Liaw, ``Synthesize battery degradation modes via a diagnostic and prognostic model,'' \emph{J. Power Sources} \textbf{219}, 204 (2012). DOI: 10.1016/j.jpowsour.2012.07.016.
\bibitem{newman2004} J. Newman, K. E. Thomas-Alyea, \emph{Electrochemical Systems}, 3rd ed., Wiley-Interscience (2004). ISBN 978-0-471-47756-3.
\bibitem{bloom2005} I. Bloom, A. N. Jansen, D. P. Abraham et al., ``Differential voltage analyses of high-power, lithium-ion cells. 1. Technique and application,'' \emph{J. Power Sources} \textbf{139}, 295 (2005). DOI: 10.1016/j.jpowsour.2004.07.021.
\end{thebibliography}

\end{document}
